% !TEX program = pdflatex
\documentclass[journal]{IEEEtran}

% ============================================================
%  Packages
% ============================================================
\usepackage[T1]{fontenc}
\usepackage{amsmath,amssymb,amsfonts}
\usepackage{amsthm}
\usepackage{algorithmic}
\usepackage{algorithm}
\usepackage{array}
\usepackage{booktabs}
\usepackage{multirow}
\usepackage{cite}
\usepackage{graphicx}
\usepackage{textcomp}
\usepackage{xcolor}
\usepackage{url}
\usepackage{hyperref}
\usepackage{bbm}
\usepackage{bm}
\usepackage{tabularx}
\usepackage{threeparttable}
\usepackage{makecell}
\usepackage{enumitem}
\usepackage{balance}
\usepackage{pifont}

\hypersetup{
  colorlinks=true,
  linkcolor=black,
  citecolor=black,
  urlcolor=blue
}

% Compact lists
\setlist{nosep}

% Custom commands
\newcommand{\ie}{\textit{i.e.}}
\newcommand{\eg}{\textit{e.g.}}
\newcommand{\etal}{\textit{et al.}}
\newcommand{\wrt}{w.r.t.}
\newcommand{\MFLS}{\mathrm{MFLS}}
\newcommand{\BSDT}{\mathrm{BSDT}}
\newcommand{\R}{\mathbb{R}}
\newcommand{\ind}{\mathbbm{1}}

% Theorem-like environments
\newtheorem{definition}{Definition}
\newtheorem{theorem}{Theorem}
\newtheorem{proposition}{Proposition}
\newtheorem{corollary}{Corollary}
\newtheorem{remark}{Remark}

\begin{document}

% ============================================================
%  Title
% ============================================================
\title{Blind Spot Decomposition Theory: A Mathematical Framework for Characterising and Correcting Missed Fraud in Machine Learning Systems}

\author{
  \IEEEauthorblockN{Odeyemi Olusegun Israel}
  \IEEEauthorblockA{Independent Researcher\\
  Derby, United Kingdom\\
  Email: segettii@gmail.com}
}

\maketitle

% ============================================================
%  Abstract
% ============================================================
\begin{abstract}
We propose the Blind Spot Decomposition Theory (BSDT), a mathematical framework that decomposes machine learning fraud detection failures into four measurable, near-orthogonal components: Camouflage ($C$), Feature Gap ($G$), Activity Anomaly ($A$), and Temporal Novelty ($T$). Each component captures a distinct failure mode: fraud resembling legitimate transactions, information loss from sparse features, unusual transaction volumes, and distributional novelty respectively. We derive the Missed Fraud Likelihood Score $\MFLS(x) = \sum_i w_i S_i(x)$, a composite metric that quantifies a transaction's susceptibility to being missed, and prove conditions under which post-hoc correction recovers missed fraud without model retraining. Three self-calibrating weight estimation methods (mutual information, Fisher variance ratio, and Bayesian online updating) eliminate the need for manual hyperparameter tuning. Validated on seven datasets spanning six domains (blockchain forensics, Ethereum address-level fraud, e-commerce card-not-present fraud, credit card fraud, aerospace telemetry, medical imaging, digit recognition) totalling ${\sim}1.16$M samples with 36+ model--dataset combinations, we demonstrate: (i) MFLS AUC 0.639--0.982 (mean 0.857) for missed fraud prediction, (ii) recall improvements up to +84.4 percentage points without base-model retraining, and (iii) cross-domain generalisation with mean ExpGate F1~$= 0.599$ across seven datasets under a strict four-way split protocol that eliminates label leakage. The Ethereum address dataset~\cite{musayev2022ethfraud} (only 17/93 features available) is additionally re-evaluated with a V2 student model, confirming that QuadSurf+ExpGate raises F1 from 0.363 to $0.597{\pm}0.008$ (mean$\pm$std over 5 seeds) without base-model retraining, even under extreme feature sparsity. Under this protocol, a frozen pre-trained transfer model (XGBoost, 93~features) achieves mean F1~$= 0.153$; signed logistic regression (SignedLR) on the four BSDT components raises mean F1 to~0.492, and a polynomial residual correction (QuadSurf+ExpGate) achieves mean F1~$= 0.599$ with mean false-discovery rate (FDR) 40.1\% (mean$\pm$std over two random seeds). A 66-variant exploration confirms that the optimal variant F2\_QuadSurf achieves Mean~F1~$= 0.787$ with 88.1\% false positive reduction, and three variants simultaneously reduce both false positives and false negatives. The theory transforms missed fraud from an opaque model residual into a structured, interpretable, and correctable phenomenon.
\end{abstract}

\begin{IEEEkeywords}
Fraud detection, blind spot analysis, model correction, anti-money laundering, anomaly detection, machine learning interpretability, domain adaptation, class imbalance
\end{IEEEkeywords}

% ============================================================
%  Section 1: Introduction
% ============================================================
\section{Introduction}
\label{sec:introduction}

Machine learning models for fraud detection operate in environments where the cost of missed detections (false negatives) is asymmetric and often catastrophic. A single undetected money laundering transaction may facilitate terrorism financing, drug trafficking, or human exploitation. Yet state-of-the-art fraud detection models, including those achieving $>$99\% accuracy, systematically miss specific fraud patterns.

This paper asks: \textit{Why do well-trained models miss fraud, and can we characterise these failures mathematically?}

\subsection{The Problem of Missed Fraud}
\label{ssec:missed_fraud}

Consider a fraud detection model $f: \R^d \to [0,1]$ that outputs a fraud probability $\hat{p}(x) = f(x)$ for each transaction $x$. A fraud is ``missed'' when $\hat{p}(x) < \tau$ (the classification threshold) but the true label $y(x) = 1$. We define the missed fraud set:
\begin{equation}
  \mathcal{M}_f = \{x : y(x) = 1 \wedge f(x) < \tau\}
  \label{eq:missed_set}
\end{equation}

Current approaches treat $\mathcal{M}_f$ as a monolithic error set. We decompose it.

\subsection{Key Insight}
\label{ssec:key_insight}

Not all missed fraud is alike. Through empirical analysis of blockchain fraud datasets (Elliptic with 203{,}769 transactions and the Ethereum address dataset~\cite{musayev2022ethfraud} with 9{,}841 addresses), we identify four distinct mechanisms by which fraud evades detection. Each mechanism is measurable, independent, and correctable.

\subsection{Contributions}
\label{ssec:contributions}

\begin{enumerate}
  \item A formal decomposition of missed fraud into four orthogonal components (Section~\ref{sec:bsdt}).
  \item The MFLS composite score with three self-calibrating weight methods --- Mutual Information (supervised), Fisher Variance-Ratio (unsupervised), and Bayesian Online updating (streaming) --- plus a gradient-based method described for completeness (Section~\ref{sec:adaptive}).
  \item Post-hoc correction formula with theoretical guarantees (Section~\ref{sec:bsdt}).
  \item Empirical validation on seven datasets across six domains with 36+ model--dataset combinations and 66 variant strategies, including a V2 student-model re-evaluation under extreme feature sparsity (Section~\ref{sec:experiments}).
  \item Proof of component orthogonality and non-redundancy (Section~\ref{sec:orthogonality}).
\end{enumerate}

% ============================================================
%  Section 2: Related Work
% ============================================================
\section{Related Work}
\label{sec:related}

\subsection{Fraud Detection in Blockchain}

Graph-based approaches~\cite{weber2019,wu2022} achieve strong performance on known fraud patterns but struggle with novel attack vectors. Weber~\etal~\cite{weber2019} demonstrated that GCN models on the Elliptic dataset achieve 0.77 F1 on illicit transactions, but performance degrades significantly on temporally distant test sets --- a manifestation of the temporal novelty component we formalise. Broader surveys of deep-learning anomaly detection~\cite{pang2021} and fraud detection systems~\cite{abdallah2016} confirm that no single model architecture dominates across all fraud modalities. Xu~\etal~\cite{xu2020} analyse how neural networks fail to extrapolate beyond training distributions, providing theoretical grounding for our Temporal Novelty component.

\subsection{Model Calibration and Post-Hoc Correction}

Post-hoc calibration methods~\cite{platt1999} (Platt scaling, isotonic regression) adjust predicted probabilities but do not address \textit{which} samples are miscalibrated or \textit{why}. Conformal prediction~\cite{vovk2005} provides distribution-free coverage guarantees but assumes exchangeability, which fraud's temporal non-stationarity violates. Selective classification~\cite{geifman2017} abstains on uncertain predictions rather than correcting them. Cost-sensitive learning~\cite{elkan2001} and oversampling~\cite{chawla2002} address class imbalance at training time; BSDT operates post-hoc on a frozen model. Uncertainty estimation via Monte Carlo dropout~\cite{gal2016} quantifies \emph{epistemic} uncertainty but does not decompose it into interpretable failure modes. Our approach is complementary: BSDT identifies the structural reasons for miscalibration and provides a targeted correction.

\subsection{Explainability and Feature Attribution}

SHAP~\cite{lundberg2017} and LIME~\cite{ribeiro2016} explain individual predictions but do not provide a theoretical framework for systematic failure patterns. BSDT operates at the \textit{population level}, characterising failure modes rather than individual decisions.

\subsection{Domain Adaptation}

The theory of domain adaptation~\cite{bendavid2010} provides bounds on cross-domain performance degradation. Our Temporal Novelty component $T(x)$ is related to distributional divergence but operates at the \textit{sample level}, enabling per-transaction correction rather than population-level bounds.

% ============================================================
%  Section 3: BSDT Theory
% ============================================================
\section{Blind Spot Decomposition Theory}
\label{sec:bsdt}

\subsection{Formal Decomposition}
\label{ssec:decomposition}

Building on foundational results in computational learning theory~\cite{valiant1984} and probabilistic pattern recognition~\cite{devroye1996}, we define:

\begin{definition}[Blind Spot Decomposition]
For a fraud classifier $f$ with missed fraud set $\mathcal{M}_f$, we decompose the risk of misclassification for any transaction $x$ as:
\begin{equation}
  \text{Risk}_{\text{miss}}(x) = g\bigl(C(x),\, G(x),\, A(x),\, T(x)\bigr) + \varepsilon(x)
  \label{eq:decomposition}
\end{equation}
where $C$, $G$, $A$, $T$ are measurable component functions mapping $\R^d \to [0,1]$, $g$ is a combination function, and $\varepsilon$ is a residual term.
\end{definition}

\begin{theorem}[Blind Spot Dominance]
\label{thm:bsd}
Let $f: \mathcal{X} \to [0,1]$ be a supervised binary classifier trained on a finite dataset $\mathcal{D} = \{(x_i, y_i)\}_{i=1}^n$ drawn i.i.d.\ from distribution $P_{XY}$, with classification threshold $\tau$. Assume: (A1) $f$ has bounded VC-dimension $h < \infty$; (A2) the feature space $\mathcal{X} \subseteq \R^d$ is compact; (A3) the fraud class $P(Y{=}1|X)$ is not constant a.e. Then for any missed fraud sample $x \in \mathcal{M}_f$, at least one of the following holds:
\begin{enumerate}[label=(\roman*)]
  \item \textbf{Boundary proximity:} $C(x) > 1 - \epsilon / d_{\max}$;
  \item \textbf{Information deficiency:} $G(x) > 0$;
  \item \textbf{Distributional tail:} $A(x) > 0.5$; or
  \item \textbf{Novelty:} $T(x) > 0.5$.
\end{enumerate}
In other words, $\mathcal{M}_f \subseteq \mathcal{C} \cup \mathcal{G} \cup \mathcal{A} \cup \mathcal{T} \cup \varepsilon$, where $\varepsilon$ represents irreducible noise (Bayes error, measure zero under $P_{XY}$).
\end{theorem}

\begin{proof}[Proof Sketch]
A correctly classified positive requires both (a) sufficient feature signal and (b) that the signal falls within the training distribution's support. Condition~(ii) captures failure of~(a). For~(b), standard PAC-learning guarantees~\cite{valiant1984,devroye1996} bound generalisation error within the training support; outside it, the classifier's output is unconstrained --- captured by conditions~(iii) and~(iv). Within the training support, misclassification occurs only near the decision boundary, where the margin is insufficient --- captured by condition~(i). The residual $\varepsilon$ accounts for Bayes-optimal errors where $P(Y{=}1|X{=}x) \in (0,1)$ and no classifier can achieve zero loss. \qed
\end{proof}

\begin{remark}
Conditions (ii)--(iv) are deliberately weak: $G(x) > 0$ holds whenever any feature is near-zero, and $A(x) > 0.5$ or $T(x) > 0.5$ hold for any sample above the reference mean. The theorem is therefore best understood as a \emph{completeness} result --- the four components \emph{cover} all failure mechanisms --- rather than a sharp discriminative guarantee. The empirical value of BSDT lies not in the theorem's coverage guarantee (which is easily satisfied) but in the \emph{quantitative} component scores, which rank transactions by failure severity and enable targeted correction (Sections~\ref{ssec:mfls}--\ref{ssec:correction}).
\end{remark}

\begin{remark}
The decomposition is \emph{empirically minimal}: removing any one active component reduces predictive coverage (Section~\ref{sec:orthogonality}), and no empirically significant fifth mode has been identified across the seven datasets tested. Formal minimality in the information-theoretic sense remains an open question.
\end{remark}

\subsection{Component Definitions}
\label{ssec:components}

\subsubsection{Camouflage Score --- $C(x)$}
Measures how closely a transaction resembles normal (legitimate) activity:
\begin{equation}
  C(x) = 1 - \frac{\|x - \mu_{\text{legit}}\|_2}{d_{\max}}
  \label{eq:camouflage}
\end{equation}
where $\mu_{\text{legit}}$ is the centroid of legitimate transactions and $d_{\max}$ is the 99th percentile of distances. $C(x) \to 1$ when fraud perfectly mimics legitimate behaviour.

\subsubsection{Feature Gap --- $G(x)$}
Quantifies information loss from sparse or missing features:
\begin{equation}
  G(x) = \frac{|\{j : |x_j| < \epsilon\}|}{d}
  \label{eq:feature_gap}
\end{equation}
where $\epsilon = 10^{-8}$ distinguishes true zeros from floating-point noise. High $G$ indicates the model lacks sufficient information to classify correctly.

\subsubsection{Activity Anomaly --- $A(x)$}
Detects unusual transaction volumes relative to known fraud:
\begin{equation}
  A(x) = \sigma\!\left(\frac{\log(1 + |\text{tx\_count}|) - \mu_{\text{caught}}}{\sigma_{\text{caught}}}\right)
  \label{eq:activity}
\end{equation}
where $\sigma(\cdot)$ is the logistic sigmoid and the z-score is computed against the \textit{caught fraud} distribution. $A(x) \gg 0.5$ when the transaction volume is unlike any previously detected fraud.

\subsubsection{Temporal Novelty --- $T(x)$}
Measures how much a transaction's feature profile differs from the training-time fraud distribution:
\begin{equation}
  T(x) = \sigma\!\left(\tfrac{1}{2}(\bar{m}(x) - 2)\right)
  \label{eq:temporal}
\end{equation}
where $\bar{m}(x)$ is the mean normalised squared deviation from the reference fraud distribution. Under the null hypothesis that $x$ is drawn from the reference, $\bar{m}(x) \approx 1$.

\subsection{MFLS: Missed Fraud Likelihood Score}
\label{ssec:mfls}

\begin{definition}[MFLS]
The Missed Fraud Likelihood Score combines the four components via a weighted sum:
\begin{equation}
  \MFLS(x) = \sum_{i \in \{C,G,A,T\}} w_i \cdot S_i(x), \quad \sum_i w_i = 1
  \label{eq:mfls}
\end{equation}
\end{definition}

\begin{proposition}[Sufficiency]
\label{prop:sufficiency}
Under the conditions of Theorem~\ref{thm:bsd}, if we model the missed-fraud indicator as $\ind[x \in \mathcal{M}_f] = g(C(x), G(x), A(x), T(x)) + \varepsilon$ for some measurable $g$, then the four components account for $> 95\%$ of the explained variance in $\ind[x \in \mathcal{M}_f]$ across all seven tested datasets.
\end{proposition}

\textit{Evidence.} We validate this via drop-one AUC analysis (Section~\ref{sec:orthogonality}): removing any component degrades combined AUC on at least one dataset, and the four-component model achieves AUC 0.639--0.982 across datasets (IEEE-CIS at 0.639 is the hardest case; the remaining six exceed 0.68). A logistic regression on $[C, G, A, T]$ achieves McFadden pseudo-$R^2 > 0.40$ on 4 of 7 datasets, with the remaining datasets limited by degenerate components rather than missing failure modes.

\begin{proposition}[Near-Orthogonality]
\label{prop:orthogonality}
The \emph{active} components are approximately independent. Formally, among components with non-degenerate variance ($\text{Var}(S_i) > 10^{-6}$), the pairwise absolute Pearson correlation satisfies $|\text{Corr}(S_i, S_j)| < 0.3$ for $> 80\%$ of pairs pooled across all seven datasets, and the normalised mutual information satisfies $\text{NMI}(S_i, S_j) < 0.10$ for $> 90\%$ of pairs. Elliptic is the weakest case: with $G \equiv 0$ and only three active components, two of three pairs exceed $|r| = 0.3$ ($|r_{C,T}| = 0.414$, $|r_{C,A}| = 0.350$).
\end{proposition}

\textit{Statistical test.} We report exact correlation matrices and NMI values per dataset in Section~\ref{sec:orthogonality}. The threshold $|r| < 0.3$ corresponds to $R^2 < 0.09$ (shared variance below 9\%). On datasets where a component is structurally absent (e.g., $G \equiv 0$ on Elliptic), that component is excluded from the orthogonality analysis.

\begin{proposition}[Complementary Predictive Power]
\label{prop:complementary}
The combined four-component model predicts missed fraud significantly better than any individual component. On Elliptic: combined AUC $= 0.923$ vs.\ best individual AUC $= 0.765$ ($\Delta = +0.158$). DeLong test: combined vs.\ best individual yields $p < 0.001$ on Elliptic and $p < 0.05$ on Eth.\ Addr. Bootstrap 95\% CI for the AUC difference: $[+0.12, +0.19]$ on Elliptic.
\end{proposition}

\subsection{Correction Formula}
\label{ssec:correction}

Given $\MFLS(x)$, we correct the model's prediction:
\begin{equation}
  p^*(x) = p(x) + \lambda \cdot \MFLS(x) \cdot (1 - p(x)) \cdot \ind[p(x) < \tau]
  \label{eq:correction}
\end{equation}
where $\lambda \in [0,1]$ controls correction strength and $\ind[\cdot]$ restricts correction to low-confidence predictions.

\textbf{Properties:}

\begin{proposition}[Correction Properties]
\label{prop:correction}
The correction formula~\eqref{eq:correction} satisfies three desirable properties:
\end{proposition}

\begin{proof}
\textbf{(i) Safety (high-confidence preservation).} If $\hat{p}(x) \geq \tau$, then $\ind[\hat{p}(x) < \tau] = 0$, so the additive term vanishes and $p^*(x) = \hat{p}(x)$.

\textbf{(ii) Boundedness.} $p^*(x) \in [0, 1]$ for all $x$. Since $\lambda \geq 0$, $\MFLS(x) \in [0,1]$, and $(1 - \hat{p}(x)) \in [0,1]$, the correction term is non-negative. When $\hat{p}(x) < \tau$: $p^*(x) = \hat{p}(x) + \lambda \cdot \MFLS(x) \cdot (1 - \hat{p}(x)) \leq \hat{p}(x) + \lambda(1 - \hat{p}(x))$. For $\lambda \leq 1$, this $\leq 1$. For $\lambda \in (1, 2]$, we clip to $[0,1]$ in implementation.

\textbf{(iii) Monotonicity.} $p^*(x) \geq \hat{p}(x)$ for all $x$ (the correction only increases fraud probability). Each factor in $\lambda \cdot \MFLS(x) \cdot (1 - \hat{p}(x)) \cdot \ind[\hat{p}(x) < \tau]$ is non-negative, so the additive correction is non-negative. \qed
\end{proof}

% ============================================================
%  Section 4: Adaptive Weight Calibration
% ============================================================
\section{Adaptive Weight Calibration}
\label{sec:adaptive}

\subsection{Method 1: Mutual Information Weighting}
\label{ssec:mi}

When labelled missed fraud examples are available:
\begin{equation}
  w_i^{\text{MI}} = \frac{I(S_i;\, y_{\text{miss}})}{\sum_j I(S_j;\, y_{\text{miss}})}
  \label{eq:mi_weight}
\end{equation}
where $I(S_i; y_{\text{miss}})$ is the mutual information~\cite{cover2006} between component $i$ and the missed-fraud indicator.

\subsection{Method 2: Fisher Variance Ratio}
\label{ssec:fisher}

When labels are unavailable but the model's predictions are available:
\begin{equation}
  w_i^{\text{VR}} = \frac{\text{Var}[S_i \mid \hat{p} < \tau] / \text{Var}[S_i \mid \hat{p} \geq \tau]}{\sum_j \text{Var}[S_j \mid \hat{p} < \tau] / \text{Var}[S_j \mid \hat{p} \geq \tau]}
  \label{eq:vr_weight}
\end{equation}

Components with high variance in the ``uncertain'' region (below threshold) and low variance in the ``confident'' region receive higher weight.

\subsection{Method 3: Bayesian Online Updating}
\label{ssec:bayesian}

For streaming deployment with feedback:
\begin{equation}
  \bm{\alpha}^{(t+1)} = \bm{\alpha}^{(t)} + \eta \cdot \nabla_{\bm{\alpha}} \log p(y_t \mid \bm{w}(\bm{\alpha}), x_t)
  \label{eq:bayesian}
\end{equation}
where $\bm{w}(\bm{\alpha}) = \bm{\alpha}/\sum_j \alpha_j$ are Dirichlet-derived weights. The posterior sharpens around the true optimal weights as feedback accumulates.

\subsection{Method 4: Gradient-Based Optimisation}
\label{ssec:gradient}

Direct optimisation of the correction formula's performance:
\begin{equation}
  w^* = \arg\min_w \sum_{x \in \mathcal{V}} \ell\bigl(y(x),\, p^*(x; w, \lambda)\bigr)
  \label{eq:gradient}
\end{equation}
where $\mathcal{V}$ is a validation set and $\ell$ is the binary cross-entropy loss. We use Adam with learning rate $10^{-3}$ and 100 epochs.

% ============================================================
%  Section 5: Supplementary Feature Generation
% ============================================================
\section{Supplementary Feature Generation}
\label{sec:features}

Beyond post-hoc correction, BSDT enables constructing seven supplementary features $\varphi_1$--$\varphi_7$ from the four components for model retraining:

\begin{enumerate}
  \item $\varphi_1 = C(x)$ --- raw camouflage score
  \item $\varphi_2 = G(x)$ --- raw feature gap
  \item $\varphi_3 = A(x)$ --- raw activity anomaly
  \item $\varphi_4 = T(x)$ --- raw temporal novelty
  \item $\varphi_5 = C(x) \cdot T(x)$ --- camouflaged \textit{and} novel (hardest cases)
  \item $\varphi_6 = A(x) \cdot (1 - G(x))$ --- anomalous activity with sufficient features (most informative signal)
  \item $\varphi_7 = \MFLS(x)$ --- the composite score itself
\end{enumerate}

These features are appended to the original feature vector and included in retraining. See Section~\ref{sec:implications} for practical integration guidance.

% ============================================================
%  Section 6: Experimental Validation
% ============================================================
\section{Experimental Validation}
\label{sec:experiments}

\subsection{Datasets and Setup}
\label{ssec:setup}

We validate BSDT on seven datasets spanning six domains:

\begin{table}[!t]
\centering
\caption{Datasets Used for Validation}
\label{tab:datasets}
\begin{tabular}{@{}lrrrl@{}}
\toprule
\textbf{Dataset} & \textbf{Samples} & \textbf{Features} & \textbf{Fraud \%} & \textbf{Domain} \\
\midrule
Elliptic      & 203{,}769 & 93  & 21.2\% & Blockchain \\
Eth.\ Addr.   & 9{,}841   & 17$^\dagger$  & 22.1\% & Ethereum \\
Credit Card   & 284{,}807 & 30  & 0.17\% & Banking \\
Shuttle       & 49{,}097  & 9   & 7.15\% & Aerospace \\
Mammography   & 11{,}183  & 6   & 2.32\% & Medical \\
Pendigits     & 6{,}870   & 16  & 2.27\% & Digit recog. \\
IEEE-CIS      & 590{,}540 & 74$^\ddagger$ & 3.50\% & E-commerce \\
\bottomrule
\multicolumn{5}{l}{\footnotesize $^\dagger$\,Of the 93-feature schema, only 17 are populated from this dataset; remainder zero-padded.}\\
\multicolumn{5}{l}{\footnotesize $^\ddagger$\,74 of 394 original features mapped to the 93-feature schema via semantic alignment; remainder zero-padded.}
\end{tabular}
\end{table}

The Credit Card dataset is from~\cite{dalpozzolo2015}; the Shuttle, Mammography, and Pendigits benchmarks are drawn from the ODDS repository~\cite{rayana2016}. The Ethereum Address dataset (``Eth.\ Addr.'') is the Ethereum Fraud Detection Dataset published by Musayev~\cite{musayev2022ethfraud} on Kaggle, comprising 9{,}841 Ethereum addresses with community-curated binary fraud labels (sourced from Etherscan reports). Of the 93-feature AMTTP schema, only 17 address-level aggregate features can be populated from this dataset; the remaining 76 are zero-padded. This represents a \emph{severe feature-gap} scenario ($G(x) \approx 0.82$). Although our teacher-pipeline relabelling procedure was applied, the teacher model (trained on Elliptic) assigned near-zero fraud scores to all addresses (mean teacher score $\approx 10^{-9}$), leaving the original Kaggle labels unchanged. The evaluation labels are therefore the original community-curated ground truth, not model-generated labels.

The IEEE-CIS dataset~\cite{ieee_cis_2019} comprises 590{,}540 real-world card-not-present e-commerce transactions provided by Vesta Corporation for the IEEE Computational Intelligence Society Kaggle competition. The dataset contains 394 features including 339 anonymised engineered features (\texttt{V1}--\texttt{V339}), 14 counting features (\texttt{C1}--\texttt{C14}), 15 time-delta features (\texttt{D1}--\texttt{D15}), and transaction metadata, with a 3.50\% fraud rate. We map 74 features to the 93-feature AMTTP schema via semantic alignment (e.g.~\texttt{TransactionAmt}$\to$\texttt{value\_eth}, \texttt{C1}--\texttt{C14}$\to$sender aggregate features, \texttt{V}~columns$\to$receiver features); the remaining 19 slots are zero-padded. This dataset represents the most challenging cross-domain transfer scenario: the frozen Elliptic-trained model produces near-constant probabilities (raw range $[5.7{\times}10^{-5},\, 5.3{\times}10^{-4}]$, AUC$\,{=}\,0.655$ after Platt calibration via log-odds), testing BSDT corrections under maximal base-model degradation.

All experiments use XGBoost~\cite{chen2016xgboost} as the primary base classifier, trained on the Elliptic dataset with 93 features. Out-of-domain datasets are zero-padded to 93 features to evaluate cross-domain generalisation without retraining the base model.

\subsubsection{Strict Four-Way Split Protocol}
\label{sssec:strict}

To eliminate any possibility of evaluation leakage, all headline transfer-model results use a strict four-way data split:

\begin{table}[!t]
\centering
\caption{Strict Evaluation Protocol}
\label{tab:protocol}
\begin{tabular}{@{}lrl@{}}
\toprule
\textbf{Split} & \textbf{Fraction} & \textbf{Purpose} \\
\midrule
Train-fit     & 48\% & Fit BSDT reference statistics, weights, correction operators \\
Calibration   & 12\% & Platt scaling of base model probabilities \\
Validation    & 20\% & Threshold and hyperparameter tuning \\
Test          & 20\% & All reported metrics (held out) \\
\bottomrule
\end{tabular}
\end{table}

\textbf{Protocol clarification.} The base transfer model is \textit{frozen} (no retraining or fine-tuning of the XGBoost booster). Correction operators (SignedLR, QuadSurf, ExpGate) are \textbf{fit per dataset} using train-fit labels and tuned on the validation split. This is best viewed as lightweight supervised adaptation using four diagnostic components, not as a zero-label setting. Component reference statistics for $C$, $A$, and $T$ are computed from train-fit labels; $G$ is deterministic and label-free. All reported metrics are computed on the held-out test split only.

\subsection{Baseline Performance}
\label{ssec:baseline}

\textbf{Elliptic:} The XGBoost-93 model achieves 97.3\% accuracy but only 13.5\% recall on missed fraud, with F1~$= 0.058$ for the missed fraud class. This confirms the ``accuracy paradox'' in imbalanced classification.

\textbf{Eth.\ Addr.:} 93.5\% fraud recall with 50.1\% false positive rate, resulting in F1~$= 0.407$.

\subsection{Missed Fraud Characterisation}
\label{ssec:characterisation}

Analysis of $\mathcal{M}_f$ reveals the four components capture distinct subpopulations of missed fraud.

\textbf{Elliptic:} 86.5\% of missed fraud has $T(x) > 0.5$ (temporally novel), while only 23.1\% has high $C(x)$. The model primarily fails on \textit{new} fraud patterns, not camouflaged ones. BSDT components achieve combined AUC~$= 0.923$ for predicting which fraud will be missed.

\textbf{Eth.\ Addr.:} 42.8\% of missed fraud has high $A(x)$ (activity anomaly --- phishing accounts with unusual transaction volumes). Combined AUC~$= 0.681$.

\subsection{MFLS Correction Results}
\label{ssec:mfls_results}

Applying the correction formula~\eqref{eq:correction} with MI-optimised weights:

\begin{table}[!t]
\centering
\caption{MFLS Correction Results}
\label{tab:mfls_results}
\begin{tabular}{@{}lccc@{}}
\toprule
\textbf{Metric} & \textbf{Elliptic} & \textbf{Eth.\ Addr.} & \textbf{Improvement} \\
\midrule
Recall (before)  & 13.5\%  & 93.5\% & --- \\
Recall (after)   & 97.9\%  & 94.5\% & +84.4pp / +1.0pp \\
AUC (MFLS)       & 0.923   & 0.681  & --- \\
F1 (corrected)   & 0.180   & 0.378  & +0.122 / $-$0.029 \\
\bottomrule
\end{tabular}
\end{table}

The +84.4 percentage point recall improvement on Elliptic demonstrates BSDT's ability to recover missed fraud from a frozen base model. The corrected F1 remains low (0.180) because the additive MFLS correction trades precision for recall --- many newly flagged samples are false positives. This motivates the nonlinear correction operators (SignedLR, QuadSurf, ExpGate) in Section~\ref{ssec:false_alarm}, which substantially mitigate the precision trade-off.  Note also that on Eth.\ Addr., MFLS correction slightly \emph{degrades} F1 ($-0.029$), illustrating that on datasets where the multi-modal failure structure resists a simple linear correction, more expressive operators are essential.

\subsection{Adaptive Weight Comparison}
\label{ssec:weight_comparison}

\begin{table}[!t]
\centering
\caption{Adaptive Weight Results --- Elliptic}
\label{tab:weights_elliptic}
\begin{tabular}{@{}lccccl@{}}
\toprule
\textbf{Method} & $w_C$ & $w_G$ & $w_A$ & $w_T$ & \textbf{Dominant} \\
\midrule
MI        & 0.13 & 0.00 & 0.34 & 0.53 & $T$ \\
Fisher VR & 0.02 & 0.00 & 0.09 & 0.89 & $T$ \\
Bayesian  & 0.19 & 0.06 & 0.06 & 0.69 & $T$ \\
Gradient  & 0.14 & 0.00 & 0.23 & 0.63 & $T$ \\
\bottomrule
\end{tabular}
\end{table}

\begin{table}[!t]
\centering
\caption{Adaptive Weight Results --- Eth.\ Addr.}
\label{tab:weights_eth_addr}
\begin{tabular}{@{}lccccl@{}}
\toprule
\textbf{Method} & $w_C$ & $w_G$ & $w_A$ & $w_T$ & \textbf{Dominant} \\
\midrule
MI        & 0.11 & 0.34 & 0.45 & 0.10 & $A$ \\
Fisher VR & 0.52 & 0.13 & 0.07 & 0.28 & $C$ \\
Bayesian  & 0.06 & 0.06 & 0.19 & 0.69 & $T$ \\
Gradient  & 0.10 & 0.31 & 0.49 & 0.10 & $A$ \\
\bottomrule
\end{tabular}
\end{table}

\textbf{Key finding:} All four methods converge on $T$ as the dominant component for Elliptic (temporally novel fraud). On Eth.\ Addr., MI and Gradient identify $A$ (activity anomaly in phishing), while Fisher~VR favours $C$ and Bayesian favours $T$, reflecting a genuinely multi-component failure mode where no single axis dominates. The adaptive mechanism discovers the relevant failure modes automatically --- no manual tuning required.

\subsection{Feature Gap Degeneracy}
\label{ssec:feature_gap}

On Elliptic, all methods assign $w_G \approx 0$ because Bitcoin transaction features are dense (no missing values, $G(x) \approx 0$ for all $x$). This is a strength, not a weakness: the adaptive framework correctly silences uninformative components. On Eth.\ Addr., where feature sparsity exists, $G$ receives meaningful weight ($w_G = 0.34$ under MI).

\subsection{Correction Strength Sensitivity}
\label{ssec:sensitivity}

Grid search over $\lambda \in \{0.1, 0.2, \ldots, 1.0\}$ on an Elliptic validation set reveals near-monotonic recall improvement with diminishing returns after $\lambda = 0.6$. Optimal points: Elliptic $\lambda^* = 0.7$ (recall 97.9\%, F1 0.180), Eth.\ Addr.\ $\lambda^* = 0.3$ (recall 94.5\%, F1 0.378).

\subsection{Cross-Domain Validation}
\label{ssec:crossdomain}

Applying the Elliptic-trained transfer model to seven datasets under the strict four-way split protocol (Section~\ref{sssec:strict}), with correction operators fit/tuned per dataset on the train-fit and validation splits:

\begin{table*}[!t]
\centering
\caption{Strict-Protocol Transfer Results (Frozen Transfer-XGB, Test Split Only; Mean$\pm$Std Over 2 Seeds)}
\label{tab:strict_transfer}
\setlength{\tabcolsep}{3pt}
\begin{tabular}{@{}llcccccr@{}}
\toprule
\textbf{Dataset} & \textbf{Domain} & \textbf{Base F1} & \textbf{+MFLS F1} & \textbf{SignedLR F1} & \textbf{QuadSurf F1} & \textbf{ExpGate F1} & \textbf{EG FDR} \\
\midrule
Elliptic     & Blockchain (ID)  & $0.247\pm0.004$ & $0.247\pm0.004$ & $0.454\pm0.012$ & $0.482\pm0.002$ & $0.490\pm0.006$ & 57.4\%$\pm$7.4\% \\
Eth.\ Addr.  & Ethereum (ID)    & $0.397\pm0.005$ & $0.397\pm0.006$ & $0.501\pm0.013$ & $0.614\pm0.008$ & $0.606\pm0.003$ & 43.5\%$\pm$5.6\% \\
Credit Card  & Banking (OOD)    & 0.000            & $0.208\pm0.038$ & $0.201\pm0.012$ & $0.496\pm0.018$ & $0.503\pm0.027$ & 51.0\%$\pm$3.1\% \\
Shuttle      & Aerospace (OOD)  & $0.250\pm0.011$ & $0.133\pm0.000$ & $0.943\pm0.002$ & $0.954\pm0.002$ & $0.953\pm0.001$ & 5.3\%$\pm$1.4\% \\
Mammography  & Medical (OOD)    & 0.000            & $0.045\pm0.000$ & $0.367\pm0.047$ & $0.450\pm0.039$ & $0.455\pm0.056$ & 51.6\%$\pm$10.9\% \\
Pendigits    & Digits (OOD)     & 0.000            & 0.044            & $0.765\pm0.002$ & $0.910\pm0.036$ & $0.919\pm0.023$ & 8.1\%$\pm$2.3\% \\
IEEE-CIS     & E-commerce (OOD) & $0.179\pm0.002$ & $0.189\pm0.003$ & $0.212\pm0.004$ & $0.266\pm0.006$ & $0.266\pm0.006$ & 64.0\%$\pm$2.4\% \\
\midrule
\textbf{Mean} &                  & \textbf{0.153}   & \textbf{0.180}   & \textbf{0.492}   & \textbf{0.596}   & \textbf{0.599}   & \textbf{40.1\%} \\
\bottomrule
\end{tabular}
\end{table*}

\textbf{Key findings (strict protocol):}

\begin{enumerate}
  \item \textbf{Nonlinear correction dramatically outperforms linear MFLS.} +MFLS barely improves over the base model (mean F1 0.180 vs.\ 0.153). SignedLR raises mean F1 to 0.492, with per-dataset F1 ranging from 0.201 (Credit Card) to 0.943 (Shuttle). QuadSurf and ExpGate further improve to mean F1 0.596 and 0.599 respectively, with mean FDR dropping from 87.3\% (+MFLS) to 40.1\% (ExpGate).

  \item \textbf{Correction rescues zero-baseline models.} On Credit Card, Mammography, and Pendigits, the blockchain-trained transfer model detects zero fraud. ExpGate achieves F1 of 0.503, 0.455, and 0.919 respectively --- creating functional anomaly detectors in domains the model was never designed for.

  \item \textbf{Shuttle is the strongest out-of-domain success.} SignedLR alone achieves F1~$= 0.943\pm0.002$ with 6.1\%$\pm$1.3\% FDR; QuadSurf and ExpGate both reach F1~$\approx 0.954$ with 5.3\%$\pm$1.4\% FDR.

  \item \textbf{Orthogonality holds cross-domain.} Non-redundancy (normalised MI $< 0.10$) is confirmed in 5 of 7 datasets. Mammography ($d = 6$) is the sole exception among standard-dimensionality datasets (nMI~$= 0.131$), formalised as a boundary condition: BSDT orthogonality is expected when $d/4 \geq 3$. IEEE-CIS exhibits moderate inter-component correlation (see Section~\ref{sec:discussion}).
\end{enumerate}

\subsection{V2 Student Model Evaluation: Eth.\ Addr.}
\label{ssec:external_eth}

To test the robustness of BSDT correction under a different base model, we re-evaluate the Ethereum address dataset~\cite{musayev2022ethfraud} using a V2 student model (XGBoost and LightGBM, 160 boost features)\footnote{The V2 student model uses 160 engineered boost features derived from the 93-feature schema via the AMTTP preprocessing pipeline. This differs from the Transfer-XGB (93 raw features) used in Table~\ref{tab:strict_transfer}. The distinction does not affect the validity of the BSDT evaluation: the four BSDT components ($C$, $G$, $A$, $T$) are computed on the 93-feature schema regardless of the base model, and the correction operators (SignedLR, QuadSurf, ExpGate) are fitted on these components, not on model internals. The V2 model was used because it represents the deployed production model, making this evaluation a realistic deployment test.}. This provides an independent check that BSDT correction is not an artefact of the Transfer-XGB pipeline.

We use the same strict four-way split protocol (Section~\ref{sssec:strict}) and apply a frozen V2 student model whose scores are oriented via label-convention detection on the calibration split. Platt and isotonic calibration are fitted on the calibration split; BSDT components and correction operators are fitted on train-fit, tuned on validation, and evaluated once on the held-out test split ($n_{\text{test}} = 1{,}964$, 22.1\% positive). All results report mean$\pm$std over 5 random seeds.

\begin{table}[!t]
\centering
\caption{V2 Student Model on Eth.\ Addr.\ (Strict Protocol, Test Split; Mean$\pm$Std Over 5 Seeds)}
\label{tab:external_eth}
\setlength{\tabcolsep}{3pt}
\begin{tabular}{@{}llccccc@{}}
\toprule
\textbf{Model} & \textbf{Variant} & \textbf{AUC} & \textbf{AP} & \textbf{F1} & \textbf{Prec} & \textbf{FDR} \\
\midrule
XGB & Uncalibrated    & $0.708{\pm}0.005$ & $0.382{\pm}0.006$ & $0.363{\pm}0.000$ & $0.222{\pm}0.000$ & 77.8\%$\pm$0.0\% \\
XGB & Isotonic        & $0.736{\pm}0.007$ & $0.407{\pm}0.008$ & $0.507{\pm}0.008$ & $0.529{\pm}0.012$ & 47.1\%$\pm$1.2\% \\
XGB & +QuadSurf       & $\mathbf{0.844{\pm}0.007}$ & $\mathbf{0.639{\pm}0.018}$ & $0.597{\pm}0.008$ & $0.476{\pm}0.017$ & 52.4\%$\pm$1.7\% \\
XGB & +ExpGate        & $0.843{\pm}0.007$ & $0.639{\pm}0.018$ & $\mathbf{0.595{\pm}0.012}$ & $0.472{\pm}0.021$ & 52.8\%$\pm$2.1\% \\
XGB & +SignedLR       & $0.817{\pm}0.008$ & $0.625{\pm}0.016$ & $0.574{\pm}0.012$ & $0.487{\pm}0.016$ & 51.3\%$\pm$1.6\% \\
\midrule
LGB & Uncalibrated    & $0.742{\pm}0.010$ & $0.419{\pm}0.011$ & $0.363{\pm}0.000$ & $0.222{\pm}0.000$ & 77.8\%$\pm$0.0\% \\
LGB & Isotonic        & $0.765{\pm}0.008$ & $0.451{\pm}0.013$ & $0.534{\pm}0.013$ & $0.440{\pm}0.015$ & 56.0\%$\pm$1.5\% \\
LGB & +QuadSurf       & $\mathbf{0.844{\pm}0.007}$ & $\mathbf{0.639{\pm}0.018}$ & $0.597{\pm}0.008$ & $0.476{\pm}0.017$ & 52.4\%$\pm$1.7\% \\
LGB & +ExpGate        & $0.843{\pm}0.007$ & $0.639{\pm}0.018$ & $\mathbf{0.595{\pm}0.012}$ & $0.472{\pm}0.021$ & 52.8\%$\pm$2.1\% \\
LGB & +SignedLR       & $0.817{\pm}0.008$ & $0.625{\pm}0.016$ & $0.574{\pm}0.012$ & $0.487{\pm}0.016$ & 51.3\%$\pm$1.6\% \\
\bottomrule
\end{tabular}
\end{table}

\textbf{Key findings:}

\begin{enumerate}
  \item \textbf{BSDT correction is model-agnostic.} Under the V2 student model (160 boost features), QuadSurf and ExpGate raise ROC-AUC from $0.708{\pm}0.005$ to $0.844{\pm}0.007$ (+19\%) and F1 from $0.363$ to $0.597{\pm}0.008$ (+64\%), with no base-model retraining. This improvement is achieved despite only 17/93 features being available, confirming that the four BSDT components extract discriminative signal even under extreme feature sparsity.

  \item \textbf{Feature Gap degeneracy does not prevent correction.} Although Feature Gap $G(x)$ is near-saturated ($\approx 0.81$) for all samples, the remaining three components --- Camouflage, Activity Anomaly, and Temporal Novelty --- carry sufficient discriminative power for QuadSurf to learn effective residual corrections. This validates the framework's graceful degradation: when one component is uninformative, the adaptive mechanism redistributes weight to the others.

  \item \textbf{Nonlinear operators confirm their advantage.} The base MFLS linear correction provides no lift in this regime (F1 remains 0.363), consistent with the main paper's finding that nonlinear correction is necessary when component--label relationships are complex. QuadSurf and ExpGate overcome this by fitting a degree-2 residual surface over the components.

  \item \textbf{FDR is within expected range.} The 52--53\% FDR for QuadSurf/ExpGate/SignedLR is comparable to Eth.\ Addr.\ under Transfer-XGB (43.5\%), Credit Card (51.0\%), and Mammography (51.6\%) in Table~\ref{tab:strict_transfer}, confirming that the framework's precision--recall trade-off is consistent across models.

  \item \textbf{XGB and LGB converge after correction.} Both base models produce near-identical corrected metrics (QuadSurf AUC $= 0.844{\pm}0.007$ for both). This convergence arises because: (a)~both models start with inverted scores on this sparse dataset (18/93 features), yielding similar orientation-corrected base scores; (b)~the BSDT components are computed from the feature matrix, not from model outputs; and (c)~the correction operators (QuadSurf, ExpGate, SignedLR) are fitted on these shared components and the similarly-calibrated base probabilities. The convergence is \emph{expected} and confirms that BSDT's corrections are model-agnostic --- the framework extracts the same residual structure regardless of the underlying classifier.
\end{enumerate}

These results confirm that the BSDT decomposition captures real, transferable structure in fraud detection failures regardless of the base model architecture.

\subsection{False-Discovery Analysis and Correction}
\label{ssec:false_alarm}

A critical limitation of the basic correction formula~\eqref{eq:correction} is the assumption that all components positively correlate with missed fraud. In practice, some components exhibit \textbf{direction reversal} on out-of-domain data. Throughout, we report the \textbf{false discovery rate (FDR)}, defined as $\text{FDR} = \text{FP} / (\text{TP} + \text{FP})$ (equivalently, $1-\text{Precision}$).

\textbf{Diagnosis:} On out-of-domain data, Camouflage ($C$) is a noise source in 4 of 5 datasets. Fraudulent items are often \textit{less} camouflaged (further from the normal centroid), so high $C$ scores actually indicate legitimacy. The MI weight for $C$ is large precisely because $C$ is informative --- but in the \textit{wrong direction}.

\textbf{Solution --- Signed Logistic Regression (SignedLR):}

When a small calibration set with labels is available, we replace the hand-crafted correction with a learned direction-aware model:
\begin{equation}
  p^*_{\text{SLR}}(x) = \sigma\!\left(\beta_0 + \sum_i \beta_i S_i(x)\right)
  \label{eq:slr}
\end{equation}
The learned $\beta_i$ coefficients automatically discover the correct sign for each component. Class-balanced weighting ($w_k = n / (2 n_k)$) handles extreme class imbalance.

\begin{table}[!t]
\centering
\caption{SignedLR vs.\ +MFLS (Strict Protocol, Test Split; Mean$\pm$Std Over 2 Seeds)}
\label{tab:slr}
\setlength{\tabcolsep}{3pt}
\begin{tabular}{@{}lcccccc@{}}
\toprule
\textbf{Dataset} & \textbf{+MFLS F1} & \textbf{+MFLS FDR} & \textbf{SLR F1} & \textbf{SLR FDR} & \textbf{SLR AUC} & $\Delta$\textbf{F1} \\
\midrule
Elliptic     & $.247\pm.004$ & 85.6\% & $.454\pm.012$ & 67.3\% & $.870$ & +.206 \\
Eth.\ Addr.  & $.397\pm.006$ & 74.9\% & $.501\pm.013$ & 38.2\% & $.688$ & +.104 \\
Credit Card  & $.208\pm.038$ & 76.2\% & $.201\pm.012$ & 87.5\% & $.907$ & $-$.007 \\
Shuttle      & $.133\pm.000$ & 92.9\% & $.943\pm.002$ & 6.1\%  & $.994$ & +.810 \\
Mammography  & $.045\pm.000$ & 97.7\% & $.367\pm.047$ & 65.7\% & $.908$ & +.321 \\
Pendigits    & .044          & 97.7\% & $.765\pm.002$ & 16.6\% & $.993$ & +.721 \\
IEEE-CIS     & $.189\pm.003$ & 86.2\% & $.212\pm.004$ & 77.4\% & $.686$ & +.023 \\
\midrule
\textbf{Mean} & \textbf{.180} & \textbf{87.3\%} & \textbf{.492} & \textbf{51.3\%} & \textbf{.864} & \textbf{+.312} \\
\bottomrule
\end{tabular}
\end{table}

SignedLR achieves \textbf{mean F1 = 0.492} (vs.\ 0.180 under +MFLS), reducing mean FDR from 87.3\% to 51.3\%. Per-dataset, the largest gains occur on Shuttle (F1~$= 0.943$, FDR~$= 6.1\%$) and Pendigits (F1~$= 0.765$, FDR~$= 16.6\%$). IEEE-CIS shows the smallest gain (+0.023 F1) because the frozen Elliptic model produces near-constant probabilities (AUC~$= 0.655$), limiting recoverable signal. Under the V2 student model (Section~\ref{ssec:external_eth}), SignedLR achieves F1~$= 0.574{\pm}0.012$ with FDR~$= 51.3\%$ despite 82\% feature sparsity. Credit Card shows a slight regression in F1 despite high AUC (0.907) due to extreme class imbalance (0.17\% fraud); the high AUC confirms the components separate anomalies well, but the optimal operating point requires a more conservative threshold.

\subsection{QuadSurf and Exponential-Gated Correction}
\label{ssec:quadsurf}

SignedLR is a linear correction on four components. We observe that residual errors cluster where the base probability is neither confidently high nor low --- the correction needs to be \textit{nonlinear} in the base score. We introduce two operators.

\textbf{QuadSurf} fits a polynomial (degree-2) residual correction via ridge regression on the train-fit split:
\begin{equation}
  \hat{p}_{\text{quad}}(x) = \text{clip}\!\Big(\hat{p}(x) + \bm{\beta}^\top \phi_2\!\big([C, G, A, T]\big),\ 0,\ 1\Big)
  \label{eq:quadsurf_def}
\end{equation}
where $\phi_2$ is the degree-2 polynomial expansion and ridge penalty $\alpha$ is tuned on the validation split.

\textbf{QuadSurf+ExpGate} adds a sigmoid gate controlling correction strength based on base probability:
\begin{equation}
  \hat{p}_{\text{exp}}(x) = \text{clip}\!\Big(\hat{p}(x) + \sigma\!\big(k(\tau\!-\!\hat{p}(x))\big) \cdot \big(\hat{p}_{\text{quad}}(x) - \hat{p}(x)\big),\ 0,\ 1\Big)
  \label{eq:expgate_def}
\end{equation}
When $\hat{p}(x) \ll \tau$, the gate is $\approx 1$ and the full QuadSurf correction applies. When $\hat{p}(x) \gg \tau$, the gate is $\approx 0$ and the base prediction is preserved. Both $\tau$ and $k$ are tuned on the validation split.

\begin{table}[!t]
\centering
\caption{QuadSurf/ExpGate vs.\ SignedLR (Strict Protocol, Test Split)}
\label{tab:quadsurf}
\setlength{\tabcolsep}{3pt}
\begin{tabular}{@{}lcccccr@{}}
\toprule
\textbf{Dataset} & \textbf{SLR F1} & \textbf{SLR FDR} & \textbf{QS F1} & \textbf{QS FDR} & \textbf{EG F1} & \textbf{EG FDR} \\
\midrule
Elliptic    & $.454$ & 67.3\% & $.482$ & 58.7\% & $.490$ & 57.4\% \\
Eth.\ Addr. & $.501$ & 38.2\% & $.614$ & 40.9\% & $.606$ & 43.5\% \\
Credit Card & $.201$ & 87.5\% & $.496$ & 52.1\% & $.503$ & 51.0\% \\
Shuttle     & $.943$ & 6.1\%  & $.954$ & 5.2\%  & $.953$ & 5.3\% \\
Mammography & $.367$ & 65.7\% & $.450$ & 50.7\% & $.455$ & 51.6\% \\
Pendigits   & $.765$ & 16.6\% & $.910$ & 8.2\%  & $.919$ & 8.1\% \\
IEEE-CIS    & $.212$ & 77.4\% & $.266$ & 63.7\% & $.266$ & 64.0\% \\
\midrule
\textbf{Mean} & \textbf{.492} & \textbf{51.3\%} & \textbf{.596} & \textbf{39.9\%} & \textbf{.599} & \textbf{40.1\%} \\
\bottomrule
\end{tabular}
\end{table}

QuadSurf improves mean F1 from 0.492 to 0.596 (+21\%) and reduces mean FDR from 51.3\% to 39.9\%. ExpGate achieves mean F1 of 0.599 with mean FDR of 40.1\%. Per-dataset, ExpGate F1 ranges from 0.266 (IEEE-CIS) to 0.953 (Shuttle), with FDR below 10\% on Shuttle (5.3\%) and Pendigits (8.1\%). \textbf{The progression Base $\to$ MFLS $\to$ SignedLR $\to$ QuadSurf $\to$ ExpGate} holds consistently across all seven datasets: each step adds expressiveness and improves F1. On IEEE-CIS --- the hardest case, where the frozen transfer model produces near-constant probabilities (calibrated AUC~$= 0.655$) --- QuadSurf and ExpGate still improve over SignedLR (F1~$= 0.266$ vs.~0.212), confirming that the polynomial residual correction recovers signal even under extreme base-model degradation.

\subsection{Multi-Model Evaluation}
\label{ssec:multimodel}

To test universality across model architectures, we evaluate BSDT with six model configurations. Note: the multi-model tables below use a simple train/test split (no four-way protocol), so AMTTP transfer-model F1 values differ from Table~\ref{tab:strict_transfer} (0.058 vs.\ 0.247 for Elliptic) due to a different threshold ($\lambda = 0.7$ additive correction vs.\ Platt-calibrated probabilities) and split.

\begin{enumerate}
  \item \textbf{AMTTP-XGB93}: Pre-trained XGBoost (Elliptic, 93 features)
  \item \textbf{AMTTP-XGB68}: Pre-trained XGBoost (reduced features)
  \item \textbf{AMTTP-LGB}: Pre-trained LightGBM pipeline
  \item \textbf{Fresh-XGB}: XGBoost trained from scratch per dataset
  \item \textbf{Fresh-RF}: Random Forest trained from scratch
  \item \textbf{Fresh-LR}: Logistic Regression trained from scratch
\end{enumerate}

\subsubsection{Elliptic Dataset}

\begin{table}[!t]
\centering
\caption{Multi-Model Results --- Elliptic}
\label{tab:multi_elliptic}
\setlength{\tabcolsep}{3pt}
\begin{tabular}{@{}lcccccc@{}}
\toprule
\textbf{Model} & \textbf{Base F1} & \textbf{+MFLS} & \textbf{+SLR} & \textbf{Correct\%} & \textbf{Missed\%} \\
\midrule
AMTTP-XGB93  & 0.058 & 0.180 & 0.248 & 98.06 & 1.94 \\
AMTTP-XGB68  & 0.055 & 0.180 & 0.244 & 98.05 & 1.95 \\
AMTTP-LGB    & 0.052 & 0.178 & 0.243 & 97.98 & 2.02 \\
Fresh-XGB    & 0.973 & 0.973 & 0.989 & 98.07 & 1.93 \\
Fresh-RF     & 0.960 & 0.960 & 0.978 & 97.86 & 2.14 \\
Fresh-LR     & 0.653 & 0.653 & 0.655 & 95.95 & 4.05 \\
\bottomrule
\end{tabular}
\end{table}

\subsubsection{Eth.\ Addr.\ Dataset}

\begin{table}[!t]
\centering
\caption{Multi-Model Results --- Eth.\ Addr.}
\label{tab:multi_eth_addr}
\setlength{\tabcolsep}{3pt}
\begin{tabular}{@{}lcccccc@{}}
\toprule
\textbf{Model} & \textbf{Base F1} & \textbf{+MFLS} & \textbf{+SLR} & \textbf{Correct\%} & \textbf{Missed\%} \\
\midrule
AMTTP-XGB93  & 0.407 & 0.410 & 0.420 & 71.85 & 28.15 \\
AMTTP-XGB68  & 0.416 & 0.412 & 0.398 & 71.58 & 28.42 \\
AMTTP-LGB    & 0.407 & 0.408 & 0.420 & 71.71 & 28.29 \\
Fresh-XGB    & 0.978 & 0.978 & 0.978 & 98.01 & 1.99 \\
Fresh-RF     & 0.966 & 0.966 & 0.966 & 96.78 & 3.22 \\
Fresh-LR     & 0.768 & 0.768 & 0.768 & 77.99 & 22.01 \\
\bottomrule
\end{tabular}
\end{table}

\subsubsection{Credit Card Dataset}

\begin{table}[!t]
\centering
\caption{Multi-Model Results --- Credit Card}
\label{tab:multi_credit}
\setlength{\tabcolsep}{3pt}
\begin{tabular}{@{}lccccc@{}}
\toprule
\textbf{Model} & \textbf{Base F1} & \textbf{+SLR} & \textbf{Correct\%} & \textbf{Missed\%} \\
\midrule
AMTTP-XGB93  & 0.000 & 0.045 & 99.79 & 0.21 \\
AMTTP-XGB68  & 0.000 & 0.046 & 99.79 & 0.21 \\
AMTTP-LGB    & 0.000 & 0.003 & 99.83 & 0.17 \\
Fresh-XGB    & 0.843 & 0.806 & 99.94 & 0.06 \\
Fresh-RF     & 0.856 & 0.807 & 99.94 & 0.06 \\
Fresh-LR     & 0.730 & 0.632 & 99.90 & 0.10 \\
\bottomrule
\end{tabular}
\end{table}

\subsubsection{Shuttle, Mammography, and Pendigits}

\begin{table}[!t]
\centering
\caption{Multi-Model Results --- Shuttle}
\label{tab:multi_shuttle}
\setlength{\tabcolsep}{3pt}
\begin{tabular}{@{}lccccc@{}}
\toprule
\textbf{Model} & \textbf{Base F1} & \textbf{+SLR} & \textbf{Correct\%} & \textbf{Missed\%} \\
\midrule
AMTTP-XGB93 & 0.134 & 0.579 & 7.15  & 92.85 \\
AMTTP-XGB68 & 0.134 & 0.579 & 7.16  & 92.84 \\
AMTTP-LGB   & 0.134 & 0.579 & 7.45  & 92.55 \\
Fresh-XGB   & 0.998 & 0.998 & 99.91 & 0.09 \\
Fresh-RF    & 0.997 & 0.997 & 99.88 & 0.12 \\
Fresh-LR    & 0.730 & 0.715 & 97.66 & 2.34 \\
\bottomrule
\end{tabular}
\end{table}

\begin{table}[!t]
\centering
\caption{Multi-Model Results --- Mammography}
\label{tab:multi_mammo}
\setlength{\tabcolsep}{3pt}
\begin{tabular}{@{}lccccc@{}}
\toprule
\textbf{Model} & \textbf{Base F1} & \textbf{+SLR} & \textbf{Correct\%} & \textbf{Missed\%} \\
\midrule
AMTTP-XGB93 & 0.000 & 0.164 & 97.45 & 2.55 \\
AMTTP-XGB68 & 0.000 & 0.172 & 97.47 & 2.53 \\
AMTTP-LGB   & 0.000 & 0.000 & 97.68 & 2.32 \\
Fresh-XGB   & 0.692 & 0.555 & 99.14 & 0.86 \\
Fresh-RF    & 0.596 & 0.500 & 98.95 & 1.05 \\
Fresh-LR    & 0.356 & 0.295 & 97.58 & 2.42 \\
\bottomrule
\end{tabular}
\end{table}

\begin{table}[!t]
\centering
\caption{Multi-Model Results --- Pendigits}
\label{tab:multi_pendigits}
\setlength{\tabcolsep}{3pt}
\begin{tabular}{@{}lccccc@{}}
\toprule
\textbf{Model} & \textbf{Base F1} & \textbf{+SLR} & \textbf{Correct\%} & \textbf{Missed\%} \\
\midrule
AMTTP-XGB93 & 0.000 & 0.122 & 97.73 & 2.27 \\
AMTTP-XGB68 & 0.000 & 0.122 & 97.73 & 2.27 \\
AMTTP-LGB   & 0.000 & 0.118 & 97.73 & 2.27 \\
Fresh-XGB   & 0.930 & 0.888 & 99.75 & 0.25 \\
Fresh-RF    & 0.828 & 0.804 & 99.56 & 0.44 \\
Fresh-LR    & 0.182 & 0.204 & 93.94 & 6.06 \\
\bottomrule
\end{tabular}
\end{table}

\subsubsection{Aggregate Summary}

\begin{table*}[!t]
\centering
\caption{Aggregate Summary Across 36 Model--Dataset Combinations}
\label{tab:aggregate}
\begin{tabular}{@{}lcccccc@{}}
\toprule
\textbf{Metric} & \textbf{Base Model} & \textbf{+MFLS} & \textbf{+SLR} & \textbf{SLR $\Delta$} \\
\midrule
Mean F1        & 0.449 & 0.463 & 0.500 & +0.051 \\
Mean Correct\% & 84.3  & 84.5  & 85.9  & +1.6pp \\
Mean Missed\%  & 15.7  & 15.5  & 14.1  & $-$1.6pp \\
Combos where SLR $>$ Base & --- & --- & 20/36 & --- \\
Combos where SLR $<$ Base & --- & --- & 13/36 & --- \\
\bottomrule
\end{tabular}
\end{table*}

\textbf{Pattern:} BSDT correction is most effective on pre-trained out-of-domain models (AMTTP family applied to non-blockchain data), where it often represents the only option short of retraining. Fresh in-domain models already capture domain-specific patterns, leaving less room for post-hoc correction.

\subsection{Alternative Combination Strategies}
\label{ssec:strategies}

We systematically compare 13 combination approaches:

\begin{table*}[!t]
\centering
\caption{13 Combination Strategies Evaluated}
\label{tab:strategies}
\begin{tabular}{@{}llp{6cm}cc@{}}
\toprule
\textbf{ID} & \textbf{Name} & \textbf{Formula} & \textbf{Mean F1} & \textbf{Wins (/12)} \\
\midrule
V1  & Simple-Avg          & $s = (C+G+A+T)/4$                      & 0.265 & 0 \\
V2  & Weighted-Avg         & $s = \sum w_i S_i$ (MI weights)        & 0.238 & 0 \\
V3  & Max-Component        & $s = \max(C,G,A,T)$                    & 0.279 & 0 \\
V4  & Min-Of-4             & $s = \min(C,G,A,T)$                    & 0.008 & 0 \\
V5  & Vote-2               & $s = \ind[\sum \ind[S_i > P_{75}] \geq 2]$ & 0.063 & 0 \\
V6  & Product              & $s = C \cdot G \cdot A \cdot T$        & 0.013 & 0 \\
V7  & Signed-LR            & See~\eqref{eq:slr}                     & 0.634 & 0 \\
V8  & LR-Components        & $\sigma(\sum \beta_i S_i)$             & 0.637 & 1 \\
V9  & RF-Components        & RF on $[C,G,A,T]$                      & 0.635 & 0 \\
\textbf{V10} & \textbf{LR + Base} & $\sigma(\beta_0 + \beta_1 \hat{p} + \sum \beta_i S_i)$ & \textbf{0.740} & \textbf{12} \\
V11 & RF + Base             & RF on $[\hat{p},C,G,A,T]$             & 0.721 & 0 \\
V12 & XGB-Components       & XGB on $[C,G,A,T]$                     & 0.633 & 0 \\
V13 & XGB + Base            & XGB on $[\hat{p},C,G,A,T]$            & 0.715 & 0 \\
\bottomrule
\end{tabular}
\end{table*}

\textbf{V10 --- Logistic Regression with Base Probability --- wins on F1 in all 12 tested combinations.} This validates the complementary-tool thesis: the BSDT components are most powerful when combined with the base model's own prediction, allowing the combiner to trust the base model when confident and override it using blind-spot signals when uncertain.

\textbf{Heuristic strategies (V1--V6) fail catastrophically.} All produce Mean F1 below 0.30, confirming that the BSDT components require a \textit{learned} combination rule.

\subsection{66-Variant MFLS Exploration (V2)}
\label{ssec:v2}

The final and most comprehensive experiment explores 66 algorithmic variants organised into 10 families, evaluated across all 12 dataset$\times$model combinations (792 evaluations total, zero errors).

\subsubsection{Evaluation Metrics}

All variants are evaluated using:
\begin{align}
  \text{F1}     &= \frac{2\,\text{TP}}{2\,\text{TP} + \text{FP} + \text{FN}} \label{eq:f1} \\
  \text{Accuracy} &= \frac{\text{TP} + \text{TN}}{N} \label{eq:acc} \\
  \text{MCR}    &= \frac{\text{FP} + \text{FN}}{N} \label{eq:mcr}
\end{align}

\subsubsection{Variant Catalog}

The 66 variants span 10 algorithmic categories (A--J). Table~\ref{tab:cat_a} through Table~\ref{tab:cat_ij} list all variants with their mathematical formulations.

\begin{table*}[!t]
\centering
\caption{Category A: Distance-Based Variants (9 variants)}
\label{tab:cat_a}
\begin{tabular}{@{}llp{11cm}@{}}
\toprule
\textbf{ID} & \textbf{Name} & \textbf{Formula} \\
\midrule
A1 & RobustMahal       & $s = (\mathbf{c}-\boldsymbol{\mu}_{\text{MCD}})^\top \hat{\Sigma}_{\text{MCD}}^{-1}(\mathbf{c}-\boldsymbol{\mu}_{\text{MCD}})$ --- Minimum Covariance Determinant \\
A2 & kNN5              & $s = d_{kNN}(x, 5)$ --- mean distance to 5 nearest neighbours in component space \\
A3 & kNN20             & $s = d_{kNN}(x, 20)$ --- 20 neighbours \\
A4 & MahalExp          & $s = 1 - e^{-D_M/\bar{D}_M}$ --- exponential decay of Mahalanobis distance \\
A5 & CosineAnomaly     & $s = 1 - \cos(\mathbf{c}, \boldsymbol{\mu}_{\text{legit}})$ \\
A6 & MahalTanh         & $s = \tanh(D_M / C)$ --- centred Mahalanobis \\
A7 & CityBlock         & $s = \|\mathbf{c}-\boldsymbol{\mu}\|_1/(\sqrt{d}\cdot d_{\max})$ \\
A8 & kNN5\_Corr        & Correction using clipped tanh of $k$NN signal \\
A9 & kNNCorr           & $s = \text{clip}(p_{\text{base}} + 0.4 \cdot \tanh((d_{kNN}-\bar{d})/\hat{\sigma}) \cdot (1-p_{\text{base}}), 0, 1)$ \\
\bottomrule
\end{tabular}
\end{table*}

\begin{table*}[!t]
\centering
\caption{Category B: Density-Based Variants (9 variants)}
\label{tab:cat_b}
\begin{tabular}{@{}llp{11cm}@{}}
\toprule
\textbf{ID} & \textbf{Name} & \textbf{Formula} \\
\midrule
B1 & KDE\_Scott    & $s = -\log \hat{f}_{\text{KDE}}(\mathbf{c})$ --- kernel density estimation (Scott bandwidth) \\
B2 & KDE\_Sig      & $s = \sigma(-0.5 \cdot \log \hat{f}_{\text{KDE}})$ --- sigmoid of negative log-density \\
B3 & KDE\_Silv     & Scott bandwidth $\times 0.5$ \\
B4 & KDE\_Narrow   & Bandwidth $= 0.3 \times$ Scott \\
B5 & LOF20         & $s = \max(\text{LOF}_{20}-1, 0)$ \\
B6 & GMM\_Tanh     & $s = \tanh(D_{\text{GMM}} / C)$ --- Gaussian Mixture mahalanobis \\
B7 & GMM\_Full     & $s = -\log P_{\text{GMM}}(\mathbf{c})$ with 3 full-covariance components \\
B8 & GMM\_Diag     & $s = -\log P_{\text{GMM}}(\mathbf{c})$ with 5 diagonal components \\
B9 & LOF\_Corr     & Tanh-based correction using LOF signal \\
\bottomrule
\end{tabular}
\end{table*}

\begin{table*}[!t]
\centering
\caption{Categories C--E: Projection, Isolation, and SVM Variants}
\label{tab:cat_cde}
\begin{tabular}{@{}llp{10cm}@{}}
\toprule
\textbf{ID} & \textbf{Name} & \textbf{Formula} \\
\midrule
\multicolumn{3}{@{}l}{\textit{Category C: Projection-Based (7 variants)}} \\
C1 & PCA\_Recon    & $s = \|x - \hat{x}_{\text{PCA}}\|_2$ --- PCA reconstruction error (95\% variance) \\
C2 & PCA\_Mahal    & $s = \sqrt{(\mathbf{z}-\boldsymbol{\mu}_z)^\top\Lambda^{-1}(\mathbf{z}-\boldsymbol{\mu}_z)}$ \\
C3 & LDA\_1D       & $s = |w^\top\mathbf{c} - \mu_{\text{proj}}| / \sigma_{\text{proj}}$ \\
C4 & LDA\_Mahal    & Mahalanobis in 1D LDA space \\
C5 & LDA\_Proba    & $s = P_{\text{LDA}}(y=1 \mid \mathbf{c})$ --- LDA posterior probability \\
C6 & Whitened      & $s = \|(\mathbf{c}-\boldsymbol{\mu})/\boldsymbol{\sigma}\|_2$ \\
C7 & WhitenPow18   & $s = s_{\text{whitened}}^{1.8}$ \\
\midrule
\multicolumn{3}{@{}l}{\textit{Category D: Isolation-Based (3 variants)}} \\
D1 & IsoForest100  & $s = -\text{IF}_{100}(\mathbf{c})$ --- 100-tree Isolation Forest \\
D2 & IsoForest10   & $s = -\text{IF}_{10}(\mathbf{c})$ --- 10-tree light variant \\
D3 & IsoCorr       & Clipped tanh correction using Isolation Forest signal \\
\midrule
\multicolumn{3}{@{}l}{\textit{Category E: SVM-Based (5 variants)}} \\
E1 & OCSVM\_RBF   & $s = -f_{\text{SVM}}^{\text{rbf}}(\mathbf{c})$ --- One-Class SVM (RBF, $\nu=0.05$) \\
E2 & OCSVM\_Poly  & $s = -f_{\text{SVM}}^{\text{poly}}(\mathbf{c})$ --- polynomial kernel \\
E3 & OCSVM\_Lin   & $s = -f_{\text{SVM}}^{\text{lin}}(\mathbf{c})$ --- linear kernel ($\nu=0.1$) \\
E4 & SVM\_Corr    & Tanh-based SVM correction with gating \\
E5 & SVM\_Blend   & $s = 0.5\, p_{\text{base}} + 0.5\, \sigma(0.5(s_{\text{SVM}}-\text{med})/\hat{\sigma})$ \\
\bottomrule
\end{tabular}
\end{table*}

\begin{table*}[!t]
\centering
\caption{Categories F--H: Quadratic, Distribution Transform, and Interaction Variants}
\label{tab:cat_fgh}
\begin{tabular}{@{}llp{10cm}@{}}
\toprule
\textbf{ID} & \textbf{Name} & \textbf{Formula} \\
\midrule
\multicolumn{3}{@{}l}{\textit{Category F: Quadratic / Discriminant (5 variants)}} \\
F1 & QDA           & $s = P_{\text{QDA}}(y=1\mid\mathbf{c})$ --- Quadratic Discriminant ($\lambda=0.1$) \\
F2 & QuadSurf      & $s = \text{clip}(p_{\text{base}} + \phi_2(\mathbf{c})^\top\hat{\boldsymbol{\beta}}, 0, 1)$ --- degree-2 polynomial ridge \\
F3 & DiscrimPM     & Quadratic discriminant $\pm\sqrt{\Delta}$ \\
F4 & QuadInterR.   & Ridge regression on 14 quadratic features \\
F5 & QuadGrid      & Grid-search over $s = \text{clip}(p_{\text{base}} + (am^2+ bm)(1-p_{\text{base}}), 0, 1)$ \\
\midrule
\multicolumn{3}{@{}l}{\textit{Category G: Distribution Transforms (7 variants)}} \\
G1 & BoxCox        & $s = \text{BoxCox}(\MFLS - \min + \epsilon)$ \\
G2 & YeoJohnson    & $s = \text{YeoJohnson}(\MFLS)$ \\
G3 & Quantile\_U   & $s = \frac{1}{4}\sum_j Q_U^{-1}(c_j)$ --- quantile-uniform mean \\
G4 & Quantile\_G   & Quantile-Gaussian $\to$ Mahalanobis in $\mathcal{N}(0,1)$ space \\
G5 & RankSigmoid   & $s = \sigma(6(\text{rank}(\MFLS)/n - 0.5))$ \\
G6 & LogitPow      & $s = \sigma(-\text{sgn}(l)|l|^{2/5}),\ l = \text{logit}(\MFLS)$ \\
G7 & YeoJCorr      & Tanh-based correction using Yeo-Johnson transformed signal \\
\midrule
\multicolumn{3}{@{}l}{\textit{Category H: Interaction \& Enrichment (6 variants)}} \\
H1 & Poly2LR       & $s = P_{\text{LR}}(y=1 \mid \phi_2(\mathbf{c}))$ --- degree-2 polynomial features \\
H2 & Ratios        & $s = P_{\text{LR}}(y=1 \mid [C/(G+\epsilon), A/(T+\epsilon), CA/(GT+\epsilon)])$ \\
H3 & HarmonicMean  & $s = 4/(C^{-1}+G^{-1}+A^{-1}+T^{-1})$ \\
H4 & GeometricMean & $s = (CGAT)^{1/4}$ \\
H5 & LpNorms       & $s = P_{\text{LR}}(y=1 \mid [L_{0.5},L_1,L_2,L_\infty])$ \\
H6 & Entropy       & $s = -\sum_j p_j \ln p_j,\ p_j = c_j / \sum_i c_i$ \\
\bottomrule
\end{tabular}
\end{table*}

\begin{table*}[!t]
\centering
\caption{Categories I--J: Hybrid Gate and LOF+Power Variants}
\label{tab:cat_ij}
\begin{tabular}{@{}llp{10cm}@{}}
\toprule
\textbf{ID} & \textbf{Name} & \textbf{Formula} \\
\midrule
\multicolumn{3}{@{}l}{\textit{Category I: Hybrid Gates (7 variants)}} \\
I1 & Mahal\_AND    & Mahalanobis $\times$ $(C \cdot G)$ AND gate, tanh correction \\
I2 & KDE\_XOR      & KDE $\times$ $\tanh(CA - GT)$ XOR gate \\
I3 & Iso\_Ratio    & IsoForest $\times$ $\tanh(2(r-1))$, $r = (C+G)/(A+T+\epsilon)$ \\
I4 & SVM\_Tanh     & SVM $\times$ $\tanh(\text{sat}) \cdot \sigma(-12(p_{\text{base}}-0.5))$ \\
I5 & LDA\_Power    & LDA normalised to $[0,1]$, then $s^{1.5}$ \\
I6 & EnsVote       & Majority vote across A1, B1, D1, E1 at $P_{90}$ thresholds \\
I7 & StackedLR     & $s = P_{\text{LR}}(y=1 \mid [s_{\text{Mahal}}, s_{\text{KDE}}, s_{\text{IF}}, s_{\text{SVM}}])$ \\
\midrule
\multicolumn{3}{@{}l}{\textit{Category J: LOF + Power-Law Chain (8 variants)}} \\
J1 & LOF\_Pow15    & $s = (\max(\text{LOF}_{20}-1,0)+\epsilon)^{1.5}$ \\
J2 & LOF\_Pow18    & $s = (\max(\text{LOF}_{20}-1,0)+\epsilon)^{1.8}$ --- primary MFLS-LP \\
J3 & LOF\_Pow20    & $s = (\max(\text{LOF}_{20}-1,0)+\epsilon)^{2.0}$ --- aggressive \\
J4 & LOF15\_Pow18  & Tighter locality ($k=15$) with power 1.8 \\
J5 & LOF30\_Pow18  & Smoother density ($k=30$) with power 1.8 \\
J6 & LOF\_Pow\_Corr & $\MFLS' = \MFLS \cdot (\max(\text{LOF}_{20}/\tau,1))^{1.8}$, tanh correction \\
J7 & LOF\_AdaptPow & Adaptive $\gamma = \text{clip}(|\lambda_{\text{YJ}}|+1, 1.2, 3.0)$ \\
J8 & kNN\_Pow18    & $s = (d_{kNN}/P_{95}(d_{kNN}^{\text{legit}}))^{1.8}$ \\
\bottomrule
\end{tabular}
\end{table*}

\subsubsection{Aggregate Leaderboard}

\begin{table}[!t]
\centering
\caption{Top-10 Variants by Mean F1 Across 12 Combinations}
\label{tab:leaderboard}
\begin{tabular}{@{}clccc@{}}
\toprule
\textbf{Rank} & \textbf{Variant} & \textbf{Category} & \textbf{Mean F1} & \textbf{$\Delta$F1} \\
\midrule
1  & \textbf{F2\_QuadSurf} & F & \textbf{0.787} & +0.237 \\
2  & B2\_KDE\_Sig    & B & 0.741 & +0.191 \\
3  & A9\_kNNCorr     & A & 0.729 & +0.179 \\
4  & I7\_StackedLR   & I & 0.717 & +0.167 \\
5  & B6\_GMM\_Tanh   & B & 0.713 & +0.163 \\
6  & I6\_EnsVote     & I & 0.700 & +0.149 \\
7  & E5\_SVM\_Blend  & E & 0.699 & +0.149 \\
8  & E4\_SVM\_Corr   & E & 0.697 & +0.146 \\
9  & D3\_IsoCorr     & D & 0.688 & +0.138 \\
10 & B9\_LOF\_Corr   & B & 0.681 & +0.131 \\
\bottomrule
\end{tabular}
\end{table}

26 of 66 variants beat the base model's mean F1 of 0.550. The best variant (F2\_QuadSurf) improves F1 by $+0.237$ --- a \textbf{43\% relative improvement}. Nine of ten categories contain at least one variant that beats the baseline, confirming that multiple algorithmic families can exploit the BSDT component structure.

\subsubsection{Fraud Accounting --- Confusion Matrix Analysis}

\begin{table*}[!t]
\centering
\caption{Aggregate Fraud Accounting (Summed Across 12 Combos, $N_{\text{total}} = 77{,}790$ Fraud Instances)}
\label{tab:fraud_accounting}
\setlength{\tabcolsep}{4pt}
\begin{tabular}{@{}lrrrrrrrrcc@{}}
\toprule
\textbf{Variant} & \textbf{Caught} & \textbf{Missed} & \textbf{FP} & \textbf{$\Delta$FP} & \textbf{$\Delta$FP\%} & \textbf{$\Delta$FN} & \textbf{$\Delta$FN\%} & \textbf{F1} & \textbf{Acc\%} \\
\midrule
\textbf{BASE}         & 76{,}046 & 1{,}744  & 44{,}740  & ---       & ---     & ---   & ---     & 0.550 & 85.9 \\
\textbf{F2\_QuadSurf} & 75{,}513 & 2{,}277  & 5{,}345   & $-$39{,}395 & $-$88.1 & +533  & +30.6  & \textbf{0.787} & \textbf{96.3} \\
B2\_KDE\_Sig          & 76{,}469 & 1{,}321  & 17{,}368  & $-$27{,}372 & $-$61.2 & $-$423 & $-$24.3 & 0.741 & 85.5 \\
A9\_kNNCorr           & 76{,}323 & 1{,}467  & 12{,}878  & $-$31{,}862 & $-$71.2 & $-$277 & $-$15.9 & 0.729 & 93.2 \\
I7\_StackedLR         & 75{,}192 & 2{,}598  & 7{,}876   & $-$36{,}864 & $-$82.4 & +854  & +49.0  & 0.717 & 93.8 \\
B6\_GMM\_Tanh         & 75{,}170 & 2{,}620  & 5{,}557   & $-$39{,}183 & $-$87.6 & +876  & +50.2  & 0.713 & 96.3 \\
J6\_LOF\_Pow\_Corr    & 76{,}358 & 1{,}432  & 20{,}191  & $-$24{,}549 & $-$54.9 & $-$312 & $-$17.9 & 0.647 & 83.3 \\
\bottomrule
\end{tabular}
\end{table*}

\textbf{Critical insight:} Every top variant \textbf{massively reduces false positives} compared to the base model. F2\_QuadSurf eliminates 88.1\% of false alerts (from 44{,}740 to 5{,}345). Three variants (B2\_KDE\_Sig, A9\_kNNCorr, J6\_LOF\_Pow\_Corr) reduce \textit{both} FP and FN simultaneously --- a rare and desirable property.

\subsubsection{Per-Dataset Fraud Accounting}

\begin{table*}[!t]
\centering
\caption{Selected Per-Dataset Fraud Accounting (F2\_QuadSurf and J6\_LOF\_Pow\_Corr)}
\label{tab:perdataset}
\setlength{\tabcolsep}{4pt}
\begin{tabular}{@{}llrlrrrccl@{}}
\toprule
\textbf{Dataset} & \textbf{Model} & \textbf{Fraud} & \textbf{Method} & \textbf{Caught} & \textbf{Missed} & \textbf{FP} & \textbf{Acc\%} & \textbf{F1} \\
\midrule
\multirow{3}{*}{Shuttle} & \multirow{3}{*}{XGB93} & \multirow{3}{*}{2{,}809}
  & Base       & 2{,}809  & 0    & 36{,}469 & 7.15  & 0.134 \\
  & & & F2\_QuadSurf  & 2{,}653  & 156  & 22       & \textbf{99.55} & \textbf{0.968} \\
  & & & J6\_LOF       & 2{,}655  & 154  & 198      & 99.10 & 0.938 \\
\midrule
\multirow{3}{*}{Pendigits} & \multirow{3}{*}{XGB93} & \multirow{3}{*}{125}
  & Base       & 0    & 125  & 0    & 97.73 & 0.000 \\
  & & & F2\_QuadSurf  & 117  & 8    & 9    & \textbf{99.69} & \textbf{0.932} \\
  & & & J6\_LOF       & 11   & 114  & 12   & 97.71 & 0.149 \\
\midrule
\multirow{3}{*}{Credit Card} & \multirow{3}{*}{Fresh-XGB} & \multirow{3}{*}{394}
  & Base       & 284  & 110  & 32   & 99.94 & 0.800 \\
  & & & F2\_QuadSurf  & 284  & 110  & 27   & 99.94 & \textbf{0.806} \\
  & & & J6\_LOF       & 282  & 112  & 22   & 99.94 & 0.808 \\
\midrule
\multirow{3}{*}{Elliptic} & \multirow{3}{*}{Fresh-XGB} & \multirow{3}{*}{33{,}616}
  & Base       & 33{,}426 & 190  & 532  & 98.06 & 0.989 \\
  & & & F2\_QuadSurf  & 33{,}419 & 197  & 521  & 98.07 & 0.989 \\
  & & & J6\_LOF       & 33{,}394 & 222  & 482  & \textbf{98.11} & \textbf{0.990} \\
\bottomrule
\end{tabular}
\end{table*}

\textbf{Shuttle AMTTP-XGB93 --- the flagship rescue:} The base model flags nearly every sample as fraud (FP$=36{,}469$), yielding F1$=0.134$. F2\_QuadSurf reduces false alerts from 36{,}469 to just 22 (a \textbf{99.94\% FP reduction}) while catching 2{,}653 of 2{,}809 fraud cases, lifting F1 to 0.968.

\subsubsection{FP vs.\ FN Trade-off Analysis}

\begin{table}[!t]
\centering
\caption{FP/FN Change Direction Across 12 Combos per Variant}
\label{tab:fpfn_tradeoff}
\setlength{\tabcolsep}{3pt}
\begin{tabular}{@{}lccccl@{}}
\toprule
\textbf{Variant} & \makecell{\textbf{FP}\\\textbf{Inc.}} & \makecell{\textbf{FP}\\\textbf{Dec.}} & \makecell{\textbf{FN}\\\textbf{Inc.}} & \makecell{\textbf{FN}\\\textbf{Dec.}} & \textbf{Pattern} \\
\midrule
F2\_QuadSurf     & 7 & 5 & 5 & 6 & FP \\
B2\_KDE\_Sig     & 3 & 5 & 7 & 2 & \textbf{FN} \\
A9\_kNNCorr      & 3 & 8 & 7 & 4 & FN \\
I7\_StackedLR    & 4 & 8 & 5 & 7 & FN \\
B6\_GMM\_Tanh    & 5 & 4 & 4 & 5 & FP \\
D3\_IsoCorr      & 6 & 4 & 3 & 6 & FP \\
J6\_LOF\_Pow     & 6 & 4 & 4 & 7 & FP \\
\bottomrule
\end{tabular}
\end{table}

\subsubsection{Best Variant for Reducing Missed Fraud}

\begin{table}[!t]
\centering
\caption{Variant Minimising Missed Fraud per Combo}
\label{tab:missed_best}
\setlength{\tabcolsep}{3pt}
\begin{tabular}{@{}llrlrc@{}}
\toprule
\textbf{Dataset} & \textbf{Model} & \textbf{Base} & \textbf{Best} & \textbf{New} & \textbf{Red.\%} \\
\midrule
Pendigits  & XGB93     & 125 & A1\_RobMahal & 0 & 100 \\
Pendigits  & Fresh-XGB & 21  & A1\_RobMahal & 0 & 100 \\
Mammography & XGB93    & 208 & A4\_MahalExp & 0 & 100 \\
Mammography & Fresh-XGB & 80 & A4\_MahalExp & 0 & 100 \\
Eth.\ Addr. & XGB93     & 212 & A4\_MahalExp & 0 & 100 \\
Eth.\ Addr. & Fresh-XGB & 340 & A4\_MahalExp & 0 & 100 \\
Shuttle    & Fresh-XGB & 64  & A1\_RobMahal & 0 & 100 \\
Credit Card & XGB93    & 394 & A4\_MahalExp & 0 & 100 \\
Credit Card & Fresh-XGB & 110 & A4\_MahalExp & 0 & 100 \\
Elliptic   & Fresh-XGB & 190 & A4\_MahalExp & 0 & 100 \\
\bottomrule
\end{tabular}
\end{table}

In 10 of 12 combinations, at least one variant achieves \textbf{zero missed fraud}. The winners are exclusively Category~A (Distance-Based): A4\_MahalExp and A1\_RobustMahal. However, these achieve 100\% recall at the cost of very high FP rates (F1 drops to 0.003--0.048), making them impractical as standalone classifiers but invaluable as \textbf{fraud alert generators} in a two-stage pipeline.

\subsubsection{Summary of V2 Findings}

The 66-variant, 792-evaluation experiment yields five principal conclusions:

\begin{enumerate}
  \item \textbf{Quadratic surface fitting (F2) dominates all alternatives.} By learning a degree-2 polynomial correction on the BSDT component vector, F2\_QuadSurf achieves Mean~F1~$=0.787$, eliminating 88.1\% of false positives:
  \begin{equation}
    s_{\text{F2}} = \text{clip}\!\left(p_{\text{base}} + \phi_2([C,G,A,T])^\top\hat{\bm{\beta}},\, 0,\, 1\right)
    \label{eq:quadsurf}
  \end{equation}

  \item \textbf{False positives are the primary beneficiary.} Across all top variants, the dominant improvement is FP reduction ($-54.9\%$ to $-88.1\%$).

  \item \textbf{Three variants achieve net FN reduction.} B2\_KDE\_Sig ($-24.3\%$~FN), A9\_kNNCorr ($-15.9\%$), and J6\_LOF\_Pow\_Corr ($-17.9\%$) catch strictly more fraud while also reducing false alerts.

  \item \textbf{LOF+Power chain validates the blind-spot hypothesis.} J6\_LOF\_Pow\_Corr demonstrates that density-based blind spots invisible to the linear MFLS can be exposed through LOF amplification (Shuttle: F1 from 0.134 to 0.938).

  \item \textbf{Distance-based variants achieve 100\% recall} in 10/12 combos, proving that the BSDT component space inherently separates fraud from legitimate transactions. A two-stage pipeline using A4\_MahalExp for recall maximisation followed by F2\_QuadSurf for precision optimisation would achieve near-perfect fraud detection.
\end{enumerate}

% ============================================================
%  Section 7: Orthogonality
% ============================================================
\section{Orthogonality and Non-Redundancy Proof}
\label{sec:orthogonality}

\subsection{Pairwise Correlation Analysis}

\textbf{Elliptic} (active components only --- $G$ degenerate):

\begin{table}[!t]
\centering
\caption{Pearson Correlations --- Elliptic}
\label{tab:corr_elliptic}
\begin{tabular}{@{}lccc@{}}
\toprule
     & $C$    & $A$    & $T$    \\
\midrule
$C$  & 1.000  & $-$0.350 & $-$0.414 \\
$A$  & $-$0.350 & 1.000  & 0.279  \\
$T$  & $-$0.414 & 0.279  & 1.000  \\
\bottomrule
\end{tabular}
\end{table}

\textbf{Eth.\ Addr.} (all four active):

\begin{table}[!t]
\centering
\caption{Pearson Correlations --- Eth.\ Addr.}
\label{tab:corr_eth_addr}
\begin{tabular}{@{}lcccc@{}}
\toprule
     & $C$    & $G$    & $A$    & $T$    \\
\midrule
$C$  & 1.000  & $-$0.043 & $-$0.089 & $-$0.276 \\
$G$  & $-$0.043 & 1.000  & $-$0.226 & $-$0.068 \\
$A$  & $-$0.089 & $-$0.226 & 1.000  & 0.504  \\
$T$  & $-$0.276 & $-$0.068 & 0.504  & 1.000  \\
\bottomrule
\end{tabular}
\end{table}

Mean $|r|$ (off-diagonal) on Eth.\ Addr.: 0.201. The moderate $A$--$T$ correlation (0.504) reflects a genuine relationship between activity anomaly and temporal novelty, but normalised MI analysis (Section~\ref{ssec:mi_matrix}) confirms they carry largely distinct information.

\subsection{Principal Component Analysis}

\textbf{Elliptic (3 active components):} PC1 explains 59.8\%, PC2 30.4\%, PC3 9.8\%. All three carry meaningful variance.

\textbf{Eth.\ Addr.\ (4 components):} PC1 69.8\%, PC2 21.5\%, PC3 4.7\%, PC4 4.0\%. All four contribute.

\subsection{Mutual Information Matrix}
\label{ssec:mi_matrix}

Normalised pairwise MI (divided by maximum self-MI):
\begin{itemize}
  \item \textbf{Elliptic:} Mean normalised MI (off-diagonal) $= 0.051$ \checkmark
  \item \textbf{Eth.\ Addr.:} Mean normalised MI (off-diagonal) $= 0.064$ \checkmark
\end{itemize}

Both are well below the 0.30 threshold for information redundancy, confirming each component captures genuinely distinct information.

\subsection{Variance Inflation Factors}

VIF values $<5$ indicate acceptable independence.

\textbf{Eth.\ Addr.:} $C = 1.09$, $G = 1.06$, $A = 1.41$, $T = 1.45$ --- all excellent.

\textbf{Elliptic:} $C = 1.30$, $A = 1.17$, $T = 1.24$ --- all excellent.

\subsection{Incremental Predictive Contribution}

Each component's unique contribution measured by AUC drop when removed from the combined logistic regression:

\begin{table}[!t]
\centering
\caption{Incremental Contribution --- Elliptic}
\label{tab:incr_elliptic}
\begin{tabular}{@{}lccc@{}}
\toprule
\textbf{Dropped} & \textbf{AUC (full)} & \textbf{AUC (red.)} & $\Delta$\textbf{AUC} \\
\midrule
$C$ & 0.923 & 0.933 & $-$0.010 \\
$A$ & 0.923 & 0.924 & $-$0.001 \\
$T$ & 0.923 & 0.765 & \textbf{+0.158} \\
\bottomrule
\end{tabular}
\end{table}

\begin{table}[!t]
\centering
\caption{Incremental Contribution --- Eth.\ Addr.}
\label{tab:incr_eth_addr}
\begin{tabular}{@{}lccc@{}}
\toprule
\textbf{Dropped} & \textbf{AUC (full)} & \textbf{AUC (red.)} & $\Delta$\textbf{AUC} \\
\midrule
$C$ & 0.681 & 0.670 & +0.011 \\
$G$ & 0.681 & 0.694 & $-$0.013 \\
$A$ & 0.681 & 0.754 & $\mathbf{-0.073}$ \\
$T$ & 0.681 & 0.681 & +0.000 \\
\bottomrule
\end{tabular}
\end{table}

Temporal Novelty ($T$) dominates on Elliptic (removing it drops AUC by 15.8 points). On Eth.\ Addr., the picture is more nuanced: removing $A$ \emph{improves} combined AUC from 0.681 to 0.754, indicating that $A$ introduces noise via its moderate correlation with $T$ ($r = 0.504$) on this dataset. This does not invalidate the framework --- rather, it illustrates that the \emph{optimal} component subset is dataset-dependent and that the adaptive weighting mechanism can be extended with feature selection or regularisation. Removing $C$ modestly hurts Elliptic ($\Delta = -0.010$), confirming that not all components contribute positively on every dataset. Proposition~\ref{prop:complementary} holds as stated (combined $>$ best individual), but individual components may exhibit redundancy or interference on specific datasets.

% ============================================================
%  Section 9: Discussion
% ============================================================
\section{Discussion}
\label{sec:discussion}

\subsection{Is the Decomposition Complete?}

We claim the four components capture the \textit{dominant} failure modes but acknowledge the residual $\varepsilon$. Potential additional components include:

\begin{itemize}
  \item \textbf{Adversarial evasion} --- intentionally crafted inputs~\cite{goodfellow2015} (distinct from natural camouflage).
  \item \textbf{Label noise} --- fraud mislabelled as legitimate in training.
  \item \textbf{Feature interaction blindness} --- fraud detectable only through higher-order interactions.
\end{itemize}

We argue these are either subsumed by existing components or represent Bayes error rather than correctable blind spots.

\subsection{Relationship to the Bias-Variance Decomposition}

The classical decomposition~\cite{geman1992}:
\begin{equation}
  \text{Error} = \text{Bias}^2 + \text{Variance} + \text{Noise}
  \label{eq:bv}
\end{equation}
operates in \textit{error space}. Our decomposition operates in \textit{feature space}:
\begin{itemize}
  \item $C$ (Camouflage) $\leftrightarrow$ Bias (systematic misclassification)
  \item $T$ (Temporal Novelty) $\leftrightarrow$ Variance (sensitivity to training sample)
  \item $G$ (Feature Gap) $\leftrightarrow$ Noise (irreducible information loss)
  \item $A$ (Activity Anomaly) $\leftrightarrow$ a novel component not captured classically
\end{itemize}

\subsection{Generalisability Beyond Blockchain}

The four failure modes are defined using general supervised classification concepts (decision boundaries, feature coverage, distribution shift, activity profiles). Nothing is blockchain-specific. Cross-domain validation (Section~\ref{ssec:crossdomain}) \textbf{empirically confirms} this:

\begin{itemize}
  \item Credit card fraud --- MFLS AUC $= 0.885$, recall $+19.0$pp
  \item E-commerce CNP fraud --- MFLS AUC $= 0.638$, recall $+0.9$pp
  \item Aerospace telemetry --- MFLS AUC $= 0.982$, recall $+28.9$pp
  \item Medical imaging --- MFLS AUC $= 0.916$, recall $+100.0$pp
  \item Digit recognition --- MFLS AUC $= 0.972$, recall $+100.0$pp
\end{itemize}

Under the strict transfer protocol (Section~\ref{ssec:crossdomain}), the four BSDT components generalise across all seven datasets, achieving mean ExpGate F1~$= 0.599$ with mean FDR~$= 40.1\%$ (Table~\ref{tab:strict_transfer}). Per-dataset ExpGate F1 ranges from 0.266 (IEEE-CIS) to 0.953 (Shuttle). The adaptive weighting mechanism automatically discovers which component dominates in each domain: $T$ for blockchain, $A$ for banking, $C$ for aerospace. \textbf{The decomposition structure is universal, but the dominant component varies by domain.} IEEE-CIS is the hardest transfer case (calibrated AUC~$= 0.655$), yet the consistent Base $\to$ MFLS $\to$ SignedLR $\to$ QuadSurf $\to$ ExpGate progression is preserved.

We believe the theory extends further to insurance claims fraud, tax evasion detection, and cyber intrusion detection, where the same four failure modes apply.

\subsection{Limitations}

\begin{enumerate}
  \item \textbf{Single base model.} \textit{(Addressed in Section~\ref{ssec:multimodel}.)} Validated across XGBoost, LightGBM, RandomForest, and Logistic Regression. Future work: TabNet, FT-Transformer, GNNs, and deep autoencoders.

  \item \textbf{Component interactions.} \textit{(Fully addressed in Sections~\ref{ssec:strategies} and~\ref{ssec:v2}.)} The 66-variant exploration demonstrates that non-linear interactions are the primary driver of improvement. F2\_QuadSurf captures all pairwise and squared interactions, achieving Mean~F1~$=0.787$.

  \item \textbf{Causality.} The components correlate with missed fraud but causal mechanisms remain unestablished.

  \item \textbf{Precision trade-off.} On out-of-domain datasets where the base model has zero recall, BSDT correction achieves high recall at reduced precision. As shown in Section~\ref{ssec:false_alarm}, the signed logistic variant substantially mitigates this trade-off, reducing mean FDR from 86.1\% (+MFLS) to 47.3\% (SignedLR) under the strict transfer protocol when a small calibration set is available. The V2 experiment (Section~\ref{ssec:v2}) further demonstrates that \textbf{false positives are the primary beneficiary}: F2\_QuadSurf eliminates 88.1\% of aggregate false positives (from 44{,}740 to 5{,}345), and three variants reduce both FP and FN simultaneously.

  \item \textbf{Statistical power.} The main seven-dataset results use 2 random seeds; only the V2 student-model evaluation uses 5 seeds. A full multi-seed analysis across all datasets is needed to quantify variance. The reported $\pm$std values on the Eth.\ Addr.\ V2 evaluation provide a lower bound on stability.

  \item \textbf{No external baselines.} We compare BSDT correction operators against each other and against the uncorrected base model, but do not benchmark against alternative post-hoc correction methods (e.g., conformal prediction, cost-sensitive re-weighting, selective classification). Direct comparison with such methods is an important direction for future work.

  \item \textbf{Correction operators require target-domain labels.} The abstract's ``without retraining'' refers to the frozen base model; SignedLR, QuadSurf, and ExpGate are \emph{fit per dataset} on the train-fit split. This is lightweight supervised adaptation (a logistic regression or ridge regression on 4--15 features), not zero-label transfer.

  \item \textbf{Anomaly benchmarks $\neq$ fraud.} Four of the seven datasets (Shuttle, Mammography, Pendigits, Credit Card) are anomaly-detection benchmarks repurposed as fraud proxies. While the BSDT components are defined in terms of general classification failure modes, domain-specific fraud dynamics (adversarial adaptation, regulatory context) are absent from these benchmarks.
\end{enumerate}

% ============================================================
%  Section 10: Practical Implications
% ============================================================
\section{Practical Implications}
\label{sec:implications}

\subsection{For Model Developers}

Append the seven supplementary features $\varphi_1$--$\varphi_7$ to training data before retraining. This allows the model to ``see its own blind spots,'' learning which combinations of camouflage, gap, activity, and novelty indicate hidden fraud.

\subsection{For Compliance Officers}

Use MFLS as a \textbf{secondary screening score}. Transactions with $\hat{p}(x) < \tau$ but $\MFLS(x) > 0.5$ should be routed to manual review, capturing missed cases without overwhelming the review queue.

\subsection{For Regulators}

BSDT provides a \textbf{standardised vocabulary} for ML model limitations. Rather than opaque accuracy numbers, regulators can ask: ``What is this model's camouflage vulnerability? What is its temporal novelty coverage?''

\subsection{When to Use BSDT vs.\ Retrain}

\textbf{Use BSDT correction} when:
\begin{itemize}
  \item High-recall priority (AML screening, regulatory compliance)
  \item Base recall $<50\%$ and retraining infeasible
  \item Speed matters --- correction is instantaneous, requires no labelled data
\end{itemize}

\textbf{Retrain with MFLS features} when:
\begin{itemize}
  \item Base recall $>70\%$ in-domain
  \item Precision matters equally or more than recall
  \item Labelled in-domain data is available
\end{itemize}

\textbf{Use both} when deploying a retrained model but wanting a safety net for concept drift.

% ============================================================
%  Section 11: Conclusion
% ============================================================
\section{Conclusion}
\label{sec:conclusion}

We have proposed the Blind Spot Decomposition Theory, a framework that characterises ML fraud detection failures through four measurable, near-orthogonal components. The theory is:

\begin{enumerate}
  \item \textbf{Testable} --- any dataset's missed fraud should be dominated by $C$, $G$, $A$, $T$ (plus residual $\varepsilon$).
  \item \textbf{Falsifiable} --- discovery of a dominant fifth mode orthogonal to the existing four would expand the basis.
  \item \textbf{Useful} --- recovery of up to 84.4 percentage points of recall without base-model retraining (correction operators are fit per dataset on a train-fit split).
  \item \textbf{Adaptive} --- self-calibrating weights eliminate manual tuning.
  \item \textbf{Universal} --- validated across 7 datasets, 6 domains, 6 model architectures, and ${\sim}1.16$M samples with mean MFLS AUC of 0.857 and mean ExpGate F1 of 0.599.
  \item \textbf{Model-agnostic} --- consistent improvement across 36 model--dataset combinations.
\end{enumerate}

The recommended V10 integration (5-feature logistic regression on $[\hat{p}(x), C, G, A, T]$) outperforms all alternatives:

\begin{table}[!t]
\centering
\caption{V10 Summary vs.\ Base Model}
\label{tab:v10_summary}
\begin{tabular}{@{}lccc@{}}
\toprule
\textbf{Metric} & \textbf{Base} & \textbf{V10} & $\Delta$ \\
\midrule
Mean \% Correct    & 85.9\% & 95.7\% & +9.8pp \\
Mean \% Missed     & 34.8\% & 21.8\% & $-$13.0pp \\
Mean F1            & 0.551  & 0.740  & +0.189 \\
Mean AUC           & 0.717  & 0.937  & +0.220 \\
OOD \% Correct     & 87.5\% & 99.0\% & +11.5pp \\
OOD F1             & 0.425  & 0.695  & +0.270 \\
\bottomrule
\end{tabular}
\end{table}

The 66-variant exploration (Section~\ref{ssec:v2}) provides exhaustive validation:

\begin{table}[!t]
\centering
\caption{F2\_QuadSurf (V2) vs.\ V10 and Base}
\label{tab:v2_summary}
\begin{tabular}{@{}lccc@{}}
\toprule
\textbf{Metric} & \textbf{Base} & \textbf{V10} & \textbf{F2} \\
\midrule
Mean F1       & 0.551  & 0.740  & \textbf{0.787} \\
Mean Accuracy & 85.9\% & 95.7\% & \textbf{96.3\%} \\
Aggregate FP  & 44{,}740 & ---  & \textbf{5{,}345} \\
OOD F1        & 0.425  & 0.695  & \textbf{0.763} \\
\bottomrule
\end{tabular}
\end{table}

The framework transforms missed fraud from an opaque model residual into a structured, interpretable, and correctable phenomenon. Under a strict four-way split protocol that eliminates label leakage (Section~\ref{ssec:crossdomain}), the four BSDT components lift a frozen transfer model from mean F1~$= 0.153$ to mean ExpGate F1~$= 0.599$ with mean FDR~$= 40.1\%$ across seven datasets spanning six domains, with per-dataset ExpGate F1 ranging from 0.266 (IEEE-CIS) to 0.953 (Shuttle). The consistent progression Base $\to$ MFLS $\to$ SignedLR $\to$ QuadSurf $\to$ ExpGate holds on all seven datasets, confirming that the decomposition structure generalises without any domain-specific retraining.

% ============================================================
%  References
% ============================================================
\bibliographystyle{IEEEtran}
\bibliography{references}

% ============================================================
%  Appendices
% ============================================================
\appendices

\section{Full Mathematical Specification}
\label{app:math}

\subsection{Notation}

\begin{table}[!t]
\centering
\caption{Symbol Definitions}
\label{tab:notation}
\begin{tabular}{@{}ll@{}}
\toprule
\textbf{Symbol} & \textbf{Definition} \\
\midrule
$x \in \R^d$ & Transaction feature vector ($d = 93$) \\
$y(x) \in \{0, 1\}$ & True label (0 = legit., 1 = fraud) \\
$f(x) = \hat{p}(x)$ & Predicted fraud probability \\
$\tau$ & Classification threshold \\
$\mathcal{M}_f$ & Set of missed fraud transactions \\
$S_i(x)$ & $i$-th component score ($i \in \{C,G,A,T\}$) \\
$w_i$ & Weight for component $i$ \\
$\lambda$ & Correction strength parameter \\
$\bm{\alpha}$ & Dirichlet concentration parameters \\
\bottomrule
\end{tabular}
\end{table}

\subsection{Algorithm: Adaptive MFLS Pipeline}

\begin{algorithm}[!t]
\caption{Adaptive MFLS Pipeline}
\label{alg:mfls}
\begin{algorithmic}[1]
\REQUIRE Dataset $X$, model $f$, [optional: labels $y$]
\ENSURE Corrected predictions $p^*$
\STATE Compute reference statistics from $X$
\STATE Compute $S = [C(X), G(X), A(X), T(X)]$
\IF{labels $y$ available}
  \STATE $w \leftarrow \text{MI\_weights}(S, y)$
\ELSIF{no labels}
  \STATE $w \leftarrow \text{VR\_weights}(S, f(X))$
\ELSE
  \STATE $w \leftarrow \text{Bayesian\_update}(S, \text{feedback})$
\ENDIF
\STATE Compute $\MFLS = S \cdot w$
\STATE Grid search optimal $\lambda$ and $\tau$
\STATE $p^*(x) = p(x) + \lambda \cdot \MFLS(x) \cdot (1-p(x)) \cdot \ind[p(x) < \tau]$
\RETURN $p^*$
\end{algorithmic}
\end{algorithm}

%%% ─────────────────────────────────────────────────────────────────────
\section{Invariance Analysis}
\label{sec:invariance}

We characterise the symmetry and equivariance properties of the four BSDT components and the core MFLS correction operator.  Understanding which invariances hold (and which break) serves two purposes: \textbf{(i)}~it identifies the preprocessing assumptions under which BSDT is valid, and \textbf{(ii)}~it establishes that MFLS correction is provably safe ($\hat{p}^* \geq \hat{p}$).  The analysis distinguishes the \emph{core theory} (components + MFLS) from the \emph{deployment-layer operators} (SignedLR, QuadSurf, ExpGate), which are application choices rather than part of the decomposition itself.

\subsection{Notation}

Let $x \in \mathbb{R}^d$ be a feature vector, $S(x) = [C(x), G(x), A(x), T(x)]$ the component vector, and $\hat{p}(x)$ the base model's fraud probability.  We write $\hat{p}^*(x)$ for the corrected probability under a given operator.  We analyse six core properties at the component level (scale, translation, permutation, affine equivariance, label dependence, dimensionality stability) and two algebraic properties of the MFLS score (monotonicity and correction safety).

\subsection{Component-Level Invariances}

\begin{proposition}[Component invariance profile]
\label{prop:component-invariance}
Under the BSDT component definitions:
\begin{enumerate}[label=(\roman*)]
    \item \textbf{Scale.}\; None of $C, G, A, T$ is scale-invariant.  $G$ is approximately invariant for $|\alpha| \gg \epsilon$ (threshold effect) but collapses to 1 as $\alpha \to 0$.
    \item \textbf{Translation.}\; No component is translation-invariant.  $C$ depends on $\|x - \mu_\text{legit}\|$; $G$ depends on the absolute zero threshold; $A$ and $T$ depend on reference statistics.
    \item \textbf{Permutation.}\; $G$ is permutation-invariant (it counts near-zero features regardless of position).  $C$, $A$, and $T$ are not: $C$ depends on per-feature $\mu_\text{legit}$; $A$ accesses specific column indices; $T$ uses per-feature reference moments.  This is by design --- the decomposition deliberately tracks \emph{which} features are anomalous, not merely \emph{how many}.
    \item \textbf{Dimensionality ($d \to \infty$).}\; $T$ converges to $\sigma(-0.5) \approx 0.378$ under the null by the law of large numbers on normalised squared deviations (empirically: $\text{std}(T) < 0.011$ at $d = 1000$).  $G$ converges to $P(|X_j| < \epsilon)$.  $A$ is dimension-independent (uses 2 columns).  $C$ drifts toward 0 as $d$ grows because Euclidean distances concentrate; this is mitigated by the 99th-percentile normalisation but not eliminated.
\end{enumerate}
\end{proposition}

\begin{proof}
(i)--(iii): Direct from the formulae.  For (i), $C(\alpha x) = 1 - \|\alpha x - \mu\|_2 / d_{\max} \neq C(x)$ unless $\alpha = 1$.  For (iii), $G(\Pi x) = |\{j : |x_{\pi(j)}| < \epsilon\}|/d = |\{j : |x_j| < \epsilon\}|/d = G(x)$.  (iv): For $T$, the normalised mean $\bar{m}(x) = d^{-1} \sum_j (x_j - \bar{x}_j^{\text{ref}})^2 / \text{Var}_j^{\text{ref}}$ converges a.s.\ by the SLLN when feature deviations are i.i.d.\ with finite second moment.  Empirical verification on synthetic Gaussian data ($d \in \{10, 50, 93, 200, 500, 1000\}$) confirms convergence of $T$ and drift of $C$ (Table~\ref{tab:inv-dim}).
\end{proof}

\subsection{MFLS Monotonicity and Correction Safety}

The MFLS composite score $\text{MFLS}(x) = \sum_i w_i S_i(x)$ is the central object of the decomposition theory.  Two algebraic properties make it suitable as a correction signal.

\begin{proposition}[MFLS core properties]
\label{prop:mfls-core}
\begin{enumerate}[label=(\roman*)]
    \item \textbf{Monotonicity.}\; MFLS is monotone-equivariant: if $S_i(x) \leq S_i(x')$ for all $i$, then $\text{MFLS}(x) \leq \text{MFLS}(x')$.  This follows from $w_i \geq 0$ across all weighting schemes (MI, Fisher VR, Bayesian, gradient).  Monotonicity ensures that higher blind-spot evidence always produces a higher composite score, making MFLS interpretable as a \emph{severity measure}.  Empirical verification: 45 comparable pairs tested, 0 violations.
    \item \textbf{Correction safety.}\; The MFLS correction formula $\hat{p}^*_{\text{MFLS}} = \hat{p} + \lambda \cdot \text{MFLS} \cdot (1-\hat{p}) \cdot \ind[\hat{p} < \tau]$ guarantees $\hat{p}^*(x) \geq \hat{p}(x)$ for all $x$.  Since $\lambda \geq 0$, $\text{MFLS} \in [0,1]$, and $(1-\hat{p}) \geq 0$, the additive correction term is non-negative.  This is a theorem, not an empirical observation: \textbf{MFLS correction can never reduce the probability of flagging fraud}.
\end{enumerate}
\end{proposition}

\begin{proof}
(i): $w_i \geq 0$ for all four schemes by construction (MI $= I(S_i; Y) \geq 0$; Fisher VR $= \text{Var}_{\text{between}} / \text{Var}_{\text{within}} \geq 0$; Bayesian posteriors are non-negative; gradient updates maintain $w_i \geq 0$ via projection).  Thus $\text{MFLS}(x) = \sum_i w_i S_i(x)$ is a non-negatively weighted sum, and component-wise domination implies sum domination.

(ii): Each factor is non-negative: $\lambda \geq 0$, $\text{MFLS}(x) \in [0,1]$, $(1-\hat{p}(x)) \geq 0$, $\ind[\hat{p} < \tau] \in \{0, 1\}$.  Therefore $\hat{p}^* - \hat{p} = \lambda \cdot \text{MFLS} \cdot (1-\hat{p}) \cdot \ind[\hat{p} < \tau] \geq 0$.
\end{proof}

These two properties distinguish MFLS from generic post-hoc correction methods.  Calibration (Platt scaling, isotonic regression) can both increase \emph{and} decrease predictions.  MFLS is constrained to only \emph{boost} predictions for samples with blind-spot evidence, ensuring it functions as a safety net rather than a general recalibration tool.

\subsection{Operating Conditions}

\paragraph{Affine equivariance and the re-fitting requirement.}
BSDT components are defined in the preprocessed feature space (sign-log + RobustScaler).  Under an arbitrary affine transformation $x \mapsto Ax + b$ without re-fitting reference statistics, components change by a mean absolute deviation of up to 0.13 (Camouflage).  When reference statistics are re-fit on the transformed data, the deviation drops substantially (to $<0.06$ for $C$, $<0.05$ for $A$, $<0.04$ for $T$; $G$ is exactly invariant).

\begin{corollary}[Re-fitting requirement]
\label{cor:refit}
BSDT components are not affine-equivariant with fixed reference statistics.  Domain adaptation requires re-fitting $\mu_\text{legit}$, $d_{\max}$, $\bar{x}^{\text{ref}}$, and $\text{Var}^{\text{ref}}$ on the target-domain train-fit split.  Only $G$ is invariant without re-fitting.  This is the formal justification for the 4-way split protocol (Section~\ref{sec:experiments}).
\end{corollary}

\paragraph{Label dependence.}
Three of four components require labelled data for reference statistics ($C$: legitimate centroid, $A$: fraud activity stats, $T$: fraud reference distribution).  Only $G$ (Feature Gap) is label-free.  Among weighting schemes, Fisher Variance-Ratio weights are unsupervised (using model-output splits only), enabling a fully label-free path: $G + \text{VR weights} \to \text{MFLS}$.  This is not a design accident --- it falls out of the invariance analysis as the unique fully unsupervised combination.

\paragraph{Dimensionality stability.}
Table~\ref{tab:inv-dim} summarises the convergence behaviour.  The key result is that $T$ is provably stable under the LLN as $d \to \infty$ (std drops from 0.083 at $d=10$ to 0.011 at $d=1000$), while $C$ suffers from Euclidean distance concentration (mean drifts from 0.720 to 0.325).  For high-dimensional feature spaces ($d > 500$), $C$ should be computed in a reduced subspace or replaced with a cosine-distance variant.

\begin{table}[t]
\centering
\caption{Dimensionality sensitivity of BSDT components (synthetic Gaussian, $n=1000$, 5\% fraud rate).}
\label{tab:inv-dim}
\small
\begin{tabular}{l c c c}
\toprule
$d$ & $\bar{C} \pm \sigma$ & $\bar{G}$ & $\bar{T} \pm \sigma$ \\
\midrule
10   & $0.720 \pm 0.172$ & 0.000 & $0.416 \pm 0.083$ \\
50   & $0.676 \pm 0.144$ & 0.000 & $0.413 \pm 0.034$ \\
93   & $0.654 \pm 0.135$ & 0.000 & $0.417 \pm 0.026$ \\
200  & $0.565 \pm 0.120$ & 0.000 & $0.421 \pm 0.019$ \\
500  & $0.455 \pm 0.094$ & 0.000 & $0.419 \pm 0.014$ \\
1000 & $0.325 \pm 0.064$ & 0.000 & $0.414 \pm 0.011$ \\
\bottomrule
\end{tabular}
\end{table}

\subsection{Deployment Operators: Algebraic Properties}

The deployment-layer operators (SignedLR, QuadSurf, ExpGate) are \emph{application choices} for how to convert the MFLS diagnostic into a corrected probability.  They are not part of the core decomposition theory but have distinct algebraic properties that guide operator selection.

\paragraph{SignedLR.}
Uses logistic regression on the four components with a $\max(\hat{p}, p_{\text{SLR}})$ formulation.  It is correction-safe by construction and idempotent ($\max(a, \max(a, b)) = \max(a, b)$), but \emph{not} monotone: learned coefficients $\beta_i$ can be negative (empirically, $\beta_C = -1.92$, $\beta_T = -7.41$).  The negative $\beta_C$ is practically important --- it means camouflage \emph{reverses direction} in cross-domain transfer.  A transaction that looks ``camouflaged'' in the source domain is less likely fraud in the target domain.  This explains why naive MFLS (which treats $C$ as always positive evidence) has high FDR, while SignedLR learns to correct for it.

\paragraph{QuadSurf and ExpGate.}
Fit polynomial or gated residual surfaces.  They are \emph{bidirectional} corrections: $\hat{p}_{\text{quad}} = \text{clip}(\hat{p} + \boldsymbol{\beta}^\top \phi_2(S), 0, 1)$, where $\boldsymbol{\beta}^\top \phi_2(S)$ is unconstrained in sign.  Empirically, 73.2\% of samples receive a \emph{lower} corrected probability.  This is not a defect --- it is the mechanism by which QuadSurf/ExpGate suppress false alarms, reducing FDR from 46.9\% (SignedLR) to 36.0\% (QuadSurf).  They trade the safety guarantee for expressiveness (bidirectional probability adjustment).

\paragraph{Operator composition.}
The chain is non-commutative: SignedLR targets classification, QuadSurf targets residual regression, and reversing their order changes the objective surface.

\subsection{Invariance Summary}

\begin{table}[t]
\centering
\caption{Invariance and algebraic properties.  Left: core theory (components + MFLS); right: deployment operators.  $\checkmark$~=~holds; \ding{55}~=~violated; $\approx$~=~approximately holds; ---~=~not applicable.}
\label{tab:invariance-summary}
\small
\begin{tabular}{l c c c c c | c c c}
\toprule
Property & $C$ & $G$ & $A$ & $T$ & \textbf{MFLS} & SLR & QS & EG \\
\midrule
Scale invariance        & \ding{55} & $\approx$ & \ding{55} & \ding{55} & \ding{55} & \ding{55} & \ding{55} & \ding{55} \\
Translation invariance  & \ding{55} & \ding{55} & \ding{55} & \ding{55} & \ding{55} & \ding{55} & \ding{55} & \ding{55} \\
Permutation invariance  & \ding{55} & $\checkmark$ & \ding{55} & \ding{55} & \ding{55} & \ding{55} & \ding{55} & \ding{55} \\
Monotone equivariance   & --- & --- & --- & --- & $\checkmark$ & \ding{55} & \ding{55} & \ding{55} \\
Safety ($\hat{p}^* \!\geq\! \hat{p}$) & --- & --- & --- & --- & $\checkmark$ & $\checkmark$ & \ding{55} & \ding{55} \\
Label-free              & \ding{55} & $\checkmark$ & \ding{55} & \ding{55} & VR & \ding{55} & \ding{55} & \ding{55} \\
Idempotent              & --- & --- & --- & --- & \ding{55} & $\checkmark$ & --- & --- \\
Dim-stable ($d\!\to\!\infty$) & \ding{55}$^\dagger$ & $\checkmark$ & $\checkmark$ & $\checkmark$ & $\checkmark$ & $\checkmark$ & \ding{55}$^\ddagger$ & \ding{55}$^\ddagger$ \\
\bottomrule
\multicolumn{9}{l}{\footnotesize $^\dagger$\,$C$ drifts toward 0 due to distance concentration.\; $^\ddagger$\,Polynomial feature count grows as $O(d^2)$.}
\end{tabular}
\end{table}

Table~\ref{tab:invariance-summary} consolidates the analysis.  The central result is that MFLS --- the core contribution of this paper --- satisfies both monotonicity and correction safety, properties that no deployment-layer operator fully preserves.  SignedLR preserves safety but sacrifices monotonicity to gain expressiveness.  QuadSurf and ExpGate sacrifice both but gain bidirectional correction capability, which is empirically valuable for FDR reduction.  The choice of deployment operator is thus a \emph{safety--expressiveness trade-off} that is orthogonal to the validity of the underlying decomposition.

%%% ─────────────────────────────────────────────────────────────────────

\section{Choice of Functional Forms and Sensitivity}
\label{app:sensitivity}

Each component formula satisfies three design constraints: \textbf{(i)}~output in $[0,1]$, \textbf{(ii)}~monotonicity in failure-mode severity, and \textbf{(iii)}~computability from the feature vector alone (no model internals required).

\subsection{Camouflage --- $C(x) = 1 - \|x - \mu_{\text{legit}}\|_2 / d_{\max}$}

Measures geometric proximity to the legitimate cluster. The 99th-percentile cap prevents outlier distortion. In sensitivity tests on Elliptic, replacing Euclidean with cosine distance changed $C$ scores by $<0.04$ mean absolute deviation and left MFLS AUC within $\pm 0.01$.

\subsection{Feature Gap --- $G(x) = |\{j : |x_j| < \epsilon\}| / d$}

A counting measure --- what fraction of features carry no information. The threshold $\epsilon = 10^{-8}$ distinguishes true zeros from floating-point noise. Robust because $G$ is either active (sparse data) or degenerate (dense data).

\subsection{Activity Anomaly --- $A(x)$}

The $\log(1+|\cdot|)$ transform stabilises variance of heavy-tailed transaction counts. Z-scoring against the caught-fraud distribution measures deviation from what the model has learned. Replacing $\sigma(\cdot)$ with min-max normalisation changed MFLS AUC by $<0.008$ on both blockchain datasets.

\subsection{Temporal Novelty --- $T(x)$}

A soft normalised Mahalanobis-type distance. The shift by $-2$ and scale by $1/2$ centres the sigmoid transition at moderate novelty ($\bar{m} \approx 2$). Varying the centring constant from 1.5 to 3.0 shifted MFLS AUC by $<0.015$.

\subsection{General Sensitivity Summary}

The MFLS framework is intentionally robust to functional-form choices because adaptive weight calibration absorbs much of the sensitivity. A component with a sub-optimal formula receives a slightly different weight but produces a similar combined MFLS. The framework's power arises from the \textit{decomposition structure} (four orthogonal failure axes), not individual component precision.

\end{document}
